\documentclass[11pt]{article}
\usepackage[T1]{fontenc}
\usepackage{lmodern,microtype}
\usepackage[margin=1in]{geometry}
\usepackage{amsmath,amssymb,amsthm,mathtools}
\usepackage{booktabs,array,longtable}
\usepackage{xcolor}
\usepackage[colorlinks=true,linkcolor=blue!45!black,citecolor=blue!45!black,urlcolor=blue!45!black]{hyperref}
\usepackage{bookmark}
\usepackage{xurl}
\allowdisplaybreaks[2]
\setlength{\emergencystretch}{2em}
\numberwithin{equation}{section}
\newtheorem{theorem}{Theorem}[section]
\newtheorem{proposition}[theorem]{Proposition}
\newtheorem{lemma}[theorem]{Lemma}
\newtheorem{corollary}[theorem]{Corollary}
\theoremstyle{definition}
\newtheorem{definition}[theorem]{Definition}
\theoremstyle{remark}
\newtheorem{remark}[theorem]{Remark}
\newtheorem{example}[theorem]{Example}
\newcommand{\dd}{\,\mathrm d}
\newcommand{\ii}{\mathrm i}
\newcommand{\R}{\mathbb R}
\newcommand{\C}{\mathbb C}
\newcommand{\N}{\mathbb N}
\DeclareMathOperator{\Res}{Res}
\DeclareMathOperator{\Log}{Log}
\title{Arithmetic residuals and backward gap continuation\\in the Newman heat flow}
\author{Edward Baker\thanks{Draft prepared for Edward Baker with OpenAI Codex LLM assistance in derivation, exposition and internal checking. The founding chat records the lead as gpt-6-astra with reasoning effort max; bounded derivation and assessment helpers were explicitly dispatched as gpt-6.1-sol/ultra. The author's intended continuation effort is ultra. The continuation lead's actual model and effort are unavailable in the canonical public-history metadata; continuation conceptual work was explicitly dispatched as gpt-6-astra/ultra, with bounded proof and assessment helpers as gpt-6.1-sol/ultra. These records describe assistance and internal review, not independent specialist validation.}}
\date{19 November 2026\\\small Analytical background, finite-cell estimates and research program}
\hypersetup{pdftitle={Arithmetic residuals and backward gap continuation in the Newman heat flow},pdfauthor={Edward Baker}}
\begin{document}
\maketitle
\begin{abstract}
A useful backward estimate for a zero gap must control the cumulative
imbalance between mutual repulsion and the complete exterior interaction.
Starting from the literal theta-source heat flow, we derive regular
pair-discriminant coordinates, a chain of dipole responses in cumulative
neighboring gaps, integer split-pole moments and a separate complex
Mellin hierarchy, together with their conditional continuation estimates.

The arithmetic input is the full heat-evolved Riemann--Siegel
approximation, whose validity is distinct from the large-height regime
where its leading terms locate individual zeros. We derive its summand
heat defect and the joint quotient error in a contour readout of gap
motion. The Polymath strip criterion, Platt--Trudgian's finite-height
input and Gomila's reported certificate identify the remaining coverage
requirements and the need for complete, correctly centered comparisons.

A bounded-domain residual formula retains both lateral boundary terms.
Recentring the relative-error logarithm disk and coupling its two relevant
Fourier coefficients reduces a complete value-to-discriminant allowance.
On a fixed arithmetic cell, the refined estimator changes a negative
sufficient lower bound to a positive one inside known coverage.
A local holomorphic correction has explicit derivative, heat-defect
and off-axis mismatch bounds. Reacquiring its divisor data improves
the complete lower bound there from $0.0054147$ to $0.1493308659$.
All these refinements are also available to the equally informed direct
consumer.

On a full natural-phase family we remove apparent correction poles and
pay derivative, mismatch and cutoff jumps. Retaining the value correction
improves the complete pointwise allowance by more than $3/40$ of the
omitted leading payment. For this consumer the separately mismatch-paid
holomorphic lift is dominated by direct retained-value evaluation.

A fixed three-prime multiplier with individual paired-ellipse radii has
complete allowance below $-1/400$ on the included lower boundary through
height $30000$. Odd-core grouping with one common full-circle dyadic
phase improves or matches this estimator, but Fourier averaging forces
the same negative bound. Following the actual dyadic phase and paying
coefficient variation gives allowance below $-1/800$ at some height in
every length-$19$ interval within a natural-cutoff chart. Independently
relaxed core rotations remain the limitation; mutual phase constraints
and signed anchor coupling remain open. No new collision-free region,
positive full-strip floor or Newman bound is claimed.
\end{abstract}
\clearpage
{\fontsize{8}{8.2}\selectfont\tableofcontents}
\clearpage
% 19 November 2026. Prepared for Edward Baker with OpenAI Codex assistance.
% Lead recorded gpt-6-astra/max. Inherited results retain their source scope.
\section{The source, the time direction and the question}
\label{sec:source}

\subsection{A distinct use of the existing hierarchy}

The present investigation begins with a concrete question. Suppose a
source theorem certifies separated zeros at a positive heat time. Can
arithmetic information about the error in their evolution preserve that
separation at an earlier time where the initial localization theorem
no longer applies? The question concerns the usefulness of a quantitative
estimate along the intervening path. A formula for the velocity at one
time, or a more accurate description of an already protected region,
is not yet that estimate.

Our starting framework is the exterior-moment manuscript
\cite{VCDCMManuscript}. That work develops exact pair dynamics,
compatible common-time approximation and direct contour enclosures.
Here we select and rederive the background needed to connect those
identities to the arithmetic approximation of Polymath
\cite{Polymath2019}. The motivating discussion
\cite{VCDCMDiscussion08,VCDCMDiscussion09,VCDCMDiscussion10,
VCDCMDiscussion11,VCDCMDiscussion12} supplies three guiding distinctions:
actual gaps versus a reference density, a joint drift versus separate
moment errors, and approximation error versus the final inequality
which protects a zero configuration. The new program is organized around
these distinctions rather than around another representation alone.

\subsection{The literal theta field}

For a source coordinate $v\ge0$ put
\begin{equation}
 \Phi(v)=\sum_{n=1}^{\infty}
 \left(2\pi^2n^4e^{9v}-3\pi n^2e^{5v}\right)
 e^{-\pi n^2e^{4v}},
 \qquad \Phi(-v)=\Phi(v).
 \label{src:phi}
\end{equation}
For real physical time $t$, and then by analytic continuation, define
\begin{equation}
 H(t,z)=\int_0^{\infty}e^{tv^2}\Phi(v)\cos(zv)\dd v
       =\frac12\int_{\R}e^{tv^2}\Phi(v)e^{\ii zv}\dd v.
 \label{src:H}
\end{equation}
The source and observation kernel remain fixed throughout this paper.
With the usual completion
\[
 \xi(s)=\tfrac12s(s-1)\pi^{-s/2}\Gamma(s/2)\zeta(s),
\]
the theta Mellin representation gives
\begin{equation}
 H(0,z)=\frac18\xi\!\left(\frac{1+\ii z}{2}\right).
 \label{src:xi}
\end{equation}
This normalization is the one used in
\cite{Polymath2019,VCDCMManuscript}. In particular Newman height $x$
is twice the corresponding absolute zeta height, not the zeta height
itself. The functional equation and evenness make either sign of $\ii z$
in \eqref{src:xi} equivalent.

\begin{proposition}[Analyticity and heat law]
\label{src:heat}
The field $H$ is jointly entire in $(t,z)$, even in $z$, and real for
real $t,z$. It obeys
\begin{equation}
 H_t=-H_{zz}.
 \label{src:heat-law}
\end{equation}
For every compact real time interval $J$, $H(t,0)$ has a strictly
positive minimum and the spatial zero divisor satisfies uniformly
\begin{equation}
 n_t(R)=O(R\log(2+R)),\qquad
 \sum_{|\rho|>R}|\rho|^{-2}
       =O\!\left(\frac{\log(2+R)}R\right).
 \label{src:count}
\end{equation}
Zeros in the sum are counted with multiplicity.
\end{proposition}

\begin{proof}
Each term of \eqref{src:phi} is positive on $v\ge0$, since
$2\pi n^2e^{4v}-3>0$. The defining series and each fixed derivative
have a bound by a polynomial in $e^v$ times
$\exp(-c e^{4v})$ for some $c>0$. On compact $(t,z)$ sets this
majorizes the integrand and its derivatives after multiplication by
$\exp(\Re(t)v^2+|\Im z|v)$. Differentiation under the integral
proves the analytic assertions and \eqref{src:heat-law}. Positivity,
continuity and compactness give the lower bound at zero.

To control spatial growth, reserve half the double-exponential decay
for integrability. The supremum of $rv-c e^{4v}$ is
$O(r\log(2+r))$; the time factor contributes only a lower-order
$O((\log(2+r))^2)$ at that maximizing scale, uniformly for $t\in J$.
Thus $\log\max_{|z|\le r}|H(t,z)|=O(1+r\log(2+r))$.
Jensen's formula, using the positive lower bound at zero and the
maximum on $|z|=2R$, proves the count. Stieltjes summation by parts
bounds the inverse-square tail by a constant times
$\int_R^\infty \log(2+s)s^{-2}\dd s$.
\end{proof}

The growth is of order at most one. Evenness therefore pairs the
canonical factors for opposite zeros, removing their linear
exponentials. Proposition~\ref{src:heat} supplies the complete tail
needed for the rational traces below. It does not locate zeros,
separate them, or impose a positive sign on an arbitrary exterior
interaction.

\subsection{Backward time is a forward heat equation}

Fix a reference time $t_*$ and put $u=t_*-t$. Then
\begin{equation}
 \partial_u H=H_{zz}.
 \label{src:backward}
\end{equation}
The words ``backward continuation'' refer to decreasing physical Newman
time. They do not mean reversing the usual positive-diffusion equation
in the real spatial coordinate. For $u\ge0$ the exact field satisfies
\begin{equation}
 H(t_*-u,z)=\frac1{\sqrt\pi}\int_{\R}
 e^{-r^2}H(t_*,z+2\sqrt u\,r)\dd r.
 \label{src:gaussian}
\end{equation}
To verify this formula, insert \eqref{src:H}, interchange integrals
using the absolute source bound, and evaluate the Gaussian transform
$\pi^{-1/2}\int e^{-r^2}e^{2\ii\sqrt u\,rv}\dd r=e^{-uv^2}$.
This is a whole-line identity. Restricting the initial data to a small
interval without paying the complement does not preserve it.

The de Bruijn--Newman constant $\Lambda$ is characterized by the
all-real-zero regime $t\ge\Lambda$; the Riemann hypothesis is equivalent
to $\Lambda\le0$. The result $\Lambda\ge0$ of Rodgers and Tao
\cite{RodgersTao2020} supplies context, not an estimate used in the local
arguments. Our aim is a useful source-specific continuation estimate;
no assertion about the value of $\Lambda$ follows from the setup.

\subsection{Notation and evidence}

\begin{center}
\begin{tabular}{>{\raggedright\arraybackslash}p{0.18\textwidth}
                >{\raggedright\arraybackslash}p{0.73\textwidth}}
\toprule
Symbol & Meaning \\
\midrule
$t,u$ & Physical Newman time and reversed time $u=t_*-t$. \\
$a,b,c,g,D$ & Pair endpoints, center, actual gap $g=b-a$, and discriminant $D=g^2/4$. \\
$Q_j,M$ & Exterior denominator $(\rho_j-c)^2-D$ and complete interaction $\sum_j Q_j^{-1}$. \\
$d,\ell$ & Height-dependent reference spacing and a separately frozen normalization scale. Neither is an actual gap by definition. \\
$p,q$ & Nonnegative integer split-pole exponents unless explicitly extended. A single complex exponent later acts on the combined denominator. \\
$S_n,T_n$ & Even and odd complete exterior moments. \\
$\mathcal E,V,K$ & Gap drift, its half-defect, and the dimensionless joint forcing. \\
$P,E,\mathcal R$ & Arithmetic reference field, error $E=H-P$, and heat defect $\mathcal R=P_t+P_{zz}$. \\
\bottomrule
\end{tabular}
\end{center}

A local identity is proved under explicit product and chart assumptions.
A continuation theorem requires inequalities along a stopped time path.
A canonical acquisition proves those inequalities for the unchanged theta
source. An improved global bound additionally requires complete coverage.
The manuscript keeps these four levels distinct. Its derivations are
internal LLM-assisted work built on the cited sources; publication or
independent specialist review is not asserted.

% 19 November 2026. Drafted for Edward Baker with OpenAI Codex assistance.
% AGC-E02; dispatched configuration gpt-6.1-sol, effort ultra.
% Dispatch metadata is recorded, not independently attested backend metadata.
\section{A counted pair and its place in the gap chain}
\label{sec:pair-chain}

We derive the local laws directly from the physical equation
$H_t=-H_{zz}$, with the normalization of Section~\ref{sec:source}.
The derivation also specifies which information is lost when a pair is
viewed in isolation. These are inherited structural identities, rederived
from \cite{VCDCMManuscript,VCDCMDiscussion08}; their validity does not
constitute a new estimate preventing a collision.

\subsection{Complete products and finite parts}

The complete-divisor statements below use the following properties on a
compact real-time interval $J$. The field is real and even in $z$,
$H(t,0)\ne0$, and its even order-at-most-one product is
\begin{equation}
 H(t,z)=H(t,0)\prod_{\{\rho,-\rho\}}
                  \left(1-\frac{z^2}{\rho(t)^2}\right).
 \label{agc:pair-product}
\end{equation}
Every zero occurrence, including multiplicity, is retained; one factor
represents two opposite occurrences. The canonical growth estimates give,
uniformly on $J$,
\begin{equation}
 n_t(R)=O_J(R\log(2+R)),\qquad
 \sum_{|\rho|>R}|\rho|^{-2}=O_J(\log(2+R)/R).
 \label{agc:pair-count}
\end{equation}
The second assertion follows from the first by partial summation.
It gives locally uniform product convergence, and Cauchy's formula gives
convergence of every fixed spatial derivative. Evenness removes the linear
exponential in the order-one Hadamard product. This restriction matters:
the divisor alone does not determine logarithmic jets of an arbitrary
entire function with an unrestricted zero-free exponential factor.

Write $B=H_z/H$. At a simple zero $x(t)$, implicit differentiation gives
\begin{equation}
 x'=\frac{H_{zz}(t,x)}{H_z(t,x)}=2U_x,
 \qquad
 U_x=\lim_{z\to x}\left(B(t,z)-\frac1{z-x}\right).
 \label{agc:pair-node}
\end{equation}
For a finite polynomial the same equation reads
$x'=2\sum_{\rho\ne x}(x-\rho)^{-1}$, by differentiating its
factorization. For the complete field the finite part in
\eqref{agc:pair-node}, or the product grouping in
\eqref{agc:pair-product}, is the definition. Separate unregularized
first-power sums need not converge absolutely.

\subsection{Symmetric coordinates through a double zero}

Suppose a fixed positively oriented contour $\Gamma$ is zero-free and
counts exactly two zeros at a reference time. On a sufficiently small
time chart this remains so. With
\[
 J_k(t)=\frac1{2\pi\ii}\oint_\Gamma z^k B(t,z)\,\dd z,
 \qquad c=\frac{J_1}{2},\qquad D=\frac{J_2}{2}-c^2,
\]
the counted quadratic and local factorization are
\begin{equation}
 Q_*(t,z)=(z-c)^2-D,\qquad H(t,z)=Q_*(t,z)A(t,z),
 \qquad R(t,z)=\frac{A_z}{A}.
 \label{agc:pair-deflation}
\end{equation}
Here $A$ is jointly holomorphic and nonzero inside the contour. The
argument-principle formulas show that $c,D$ are holomorphic even when
the counted zeros merge. For a conjugation-invariant pair at real time,
$c,D$ are real. Write $z_\pm=c\pm s$, $s^2=D$, on a separated local
chart. Then $D>0$ describes two real roots, $D<0$ a conjugate pair,
and $D=0$ a double root. On a real adjacent edge the gap is
$g=z_+-z_-=2s>0$, so $D=g^2/4$.

\begin{proposition}[Pair law with the complete exterior]
\label{agc:pair-law-prop}
Let $R[a,b]=(R(b)-R(a))/(b-a)$, continued as $R_z(a)$ at
$a=b$, and set $M=-R[c-s,c+s]$. Then
\begin{equation}
 c'=R(t,c-s)+R(t,c+s),\qquad
 D'=2-4DM.
 \label{agc:pair-law}
\end{equation}
Under \eqref{agc:pair-product}--\eqref{agc:pair-count},
\begin{equation}
 M=\sum_{\rho\ \mathrm{exterior}}
       \frac1{(\rho-c)^2-D}.
 \label{agc:pair-M}
\end{equation}
This sum is absolute, includes every exterior occurrence, and continues
through exterior multiplicities and a double target while the contour
remains cleared. At a target double, $M=-R_z(t,c)$ and $D'=2$.
\end{proposition}

\begin{proof}
At a simple target, differentiating $H=Q_*A$ twice in space gives
\[
 z_+'=\frac1s+2R(t,c+s),\qquad
 z_-'=-\frac1s+2R(t,c-s).
\]
Taking the mean gives $c'$. Since $D=s^2$,
$D'=s(z_+'-z_-')=2+2s(R(c+s)-R(c-s))$, which is
\eqref{agc:pair-law}. These symmetric expressions extend holomorphically
to $s=0$ by the counted coordinates.
For an exterior occurrence, its logarithmic pole contributes
\[
 -\frac1{2s}\left(\frac1{c+s-\rho}-\frac1{c-s-\rho}\right)
       =\frac1{(\rho-c)^2-D}.
\]
The product regularizers are constant in space and cancel in this
difference. The resulting summand is $O(|\rho|^{-2})$ at infinity,
so \eqref{agc:pair-count} justifies summation and passage to the limit.
Alternatively, for $Q_*=(z-c)^2-D$ the same expression is
\begin{equation}
 M=-\frac1{2\pi\ii}\oint_\Gamma Q_*^{-1}B\,\dd z.
 \label{agc:pair-M-contour}
\end{equation}
For each exterior pole the integral is minus its rational weight;
the target contribution has total residue zero because all its poles
are inside $\Gamma$ and it decays as $z^{-3}$ at infinity.
The fixed contour proves holomorphy at all allowed multiplicities.
\end{proof}

The analytic exterior factor is invisible to polynomial-weighted count
integrals: $\oint_\Gamma z^kR\,\dd z=0$ for $k\ge0$.
Nevertheless its values at the targets enter their velocities through
the product rule. Identical current target positions can therefore have
different subsequent motion when their complete exteriors differ.
This is a local statement about deflation, not permission to change
the canonical field by an arbitrary analytic multiplier.

The first spatial variation of this analytic field can be retained
without discarding its effect on the clock. Put $Z=z-c$ and
$b=R_z(t,c)$. On a cleared contour enclosing the pair and its center,
Cauchy's formula and
$1/(Z^2-D)=1/Z^2+D/[Z^2(Z^2-D)]$ give
\begin{equation}
 \begin{split}
 D'&=2+4bD+E_D,\\
 b&=\frac1{2\pi\ii}\oint_\Gamma\frac{B}{Z^2}\,\dd z,\qquad
 E_D=\frac{4D^2}{2\pi\ii}\oint_\Gamma
                         \frac{B}{Z^2(Z^2-D)}\,\dd z.
 \end{split}
 \label{agc:pair-affine}
\end{equation}
Replacing $R$ by $B$ is valid since the added target rational functions
have all poles inside and total residue zero. If $|R_{zzz}|\le C_3$
on both center-to-root segments, Taylor's formula with integral
remainder gives $|E_D|\le(2/3)C_3|D|^2$.
The constant exterior response affects $c'$, while the affine slope
$b$ affects $D'$ already at order $D$. A small $E_D$ alone supplies
no bound on the uncertainty in $b$; the full rational kernel is exactly
$4D/Z^2+4D^2/[Z^2(Z^2-D)]=4D/(Z^2-D)$.

\subsection{Cumulative distances and the dipole response}

Now assume the \emph{complete} divisor is real and simple on an interval,
and label its ordered roots $x_n$, $n\in\mathbb Z$. Define
$g_n=x_{n+1}-x_n>0$, $c_n=(x_n+x_{n+1})/2$,
$D_n=g_n^2/4$, and let $M_n$ be \eqref{agc:pair-M} for this edge.
For $m\ge1$, put
\[
 L_{n,m}=\sum_{j=1}^m g_{n-j},\qquad
 R_{n,m}=\sum_{j=1}^m g_{n+j}.
\]
The left exterior root $x_{n-m}$ is at distances $-L_{n,m}$ and
$-(L_{n,m}+g_n)$ from the two targets. The right root
$x_{n+1+m}$ is at distances $R_{n,m}+g_n$ and $R_{n,m}$.
They exhaust the exterior, so absolute convergence proves
\begin{equation}
 \begin{split}
 M_n&=\sum_{m\ge1}\left\{
 \frac1{L_{n,m}(L_{n,m}+g_n)}+
 \frac1{R_{n,m}(R_{n,m}+g_n)}\right\},\\
 g_n'&=\frac4{g_n}-2g_nM_n.
 \end{split}
 \label{agc:pair-chain-law}
\end{equation}
This is a nonlocal equation: a gap enters many rows through cumulative
distances. It determines relative motion, leaving the common translation
coordinate to the node or center law. If only a finite real block is
known, only its actual terms have this positive-gap form. Every omitted
real or nonreal occurrence must remain in the complete interaction.

For targets $a,b=a+g$ and an exterior $\rho$,
\begin{equation}
 \frac1{(\rho-b)(\rho-a)}
   =\frac1g\left(\frac1{\rho-b}-\frac1{\rho-a}\right),\qquad
 M=-\frac{R(b)-R(a)}g=-\frac1g\int_a^b R_z(z)\,\dd z.
 \label{agc:pair-dipole}
\end{equation}
Thus the dipole interpretation is an exact difference of reciprocal
fields. At $w=\rho-(a+b)/2$ it is
$w^{-2}(1+O(g^2/w^2))$ as $|w|/g\to\infty$; the common far
response cancels before summation. $M$ is the complete response per
unit separation, not a population average.

There is also an exact spatial cancellation across edges. From
$x_n'=2U_n$ and \eqref{agc:pair-chain-law},
\begin{equation}
 \begin{split}
 g_nM_n&=\frac2{g_n}-(U_{n+1}-U_n),\\
 \sum_{n=a}^{b}g_nM_n
   &=2\sum_{n=a}^{b}\frac1{g_n}-(U_{b+1}-U_a).
 \end{split}
 \label{agc:pair-block}
\end{equation}
The proof is subtraction of the two node velocities followed by finite
telescoping. Shared interior force errors cancel if the same node
approximation is used on both edges. The identity controls a block
aggregate; it does not stop that aggregate from concentrating into a
single small gap. At a merger, the regular chart of the merging pair
may survive, but its neighboring edge charts lose exterior separation
because a formerly exterior occurrence meets a target.

\subsection{Changing density and local fluctuations}

The canonical chain has a varying logarithmic mean scale
\cite{VCDCMDiscussion09,VCDCMDiscussion10}. In the normalization
$x=2\gamma$ at $t=0$, the smooth counting main term and its derivative are
\begin{equation}
 \overline N(X)=\frac{X}{4\pi}\left(\log\frac{X}{4\pi}-1\right),
 \qquad \nu(X)=\frac1{4\pi}\log\frac{X}{4\pi},\qquad
 d_X=\nu(X)^{-1}.
 \label{agc:pair-density}
\end{equation}
Here $\nu$ differentiates the smooth main term, not the discrete
counting staircase or its error. The complete counting law alone
does not make the divisor real or give each real gap a lower bound.

Freeze a positive reference length $d$ near a chosen height and put
$\alpha_n=g_n/d$, $\mu_n=d^2M_n$, and
$\tau=(t-t_*)/d^2$. Rescaling the cumulative distances in
\eqref{agc:pair-chain-law} yields its exact dimensionless version, and
\begin{equation}
 D_n'=2-\alpha_n^2\mu_n,\qquad
 \frac{\dd\alpha_n}{\dd\tau}=\frac4{\alpha_n}-2\alpha_n\mu_n.
 \label{agc:pair-unfolded}
\end{equation}
Both the target fluctuation $\alpha_n$ and the environmental response
$\mu_n$ matter. A bound $d^2M_n\simeq2$ does not imply near cancellation
unless $g_n/d$ is controlled too. If $d$ instead varies with time,
$Y=D/d(t)^2$ obeys $Y'=D'/d^2-2Yd'/d$.
The logarithmic scale therefore cannot be inserted as a moving unit
without its derivative payment.

In the existing positive-time canonical acquisition, individual zeros
are first localized near a nonuniform clock; that input already gives
the target gap a positive lower bound. Only then are the nearby terms
compared with reference cumulative distances and the complete distant
tail paid to obtain $d^2M\simeq2$ \cite{VCDCMDiscussion10}.
The downstream interaction estimate is therefore not a derivation of
individual separation from smooth counting density alone. Extending
its domain requires additional canonical information.

As a normalization control only, a uniform infinite chain with gap $g$
has $M=2g^{-2}\sum_{m\ge1}[m(m+1)]^{-1}=2/g^2$ and $g'=0$.
Its telescoping sum checks the factors in the physical law; it is not
a canonical model of the height-dependent fluctuations.
Even the isolated quadratic heat field $z^2-2(t-T)$ obeys
\eqref{agc:pair-law}, with $M=0$, and has a double at $t=T$.
Such controls fail the literal theta-source premise, not this local
algebra. The next section organizes the complete exterior response;
the source estimate needed for continuation remains additional input.

% 19 November 2026. Drafted for Edward Baker with OpenAI Codex assistance.
% AGC-E02; dispatched configuration gpt-6.1-sol, effort ultra.
% Dispatch metadata is recorded, not independently attested backend metadata.
\section{Integer pole powers and a single complex exponent}
\label{sec:hierarchy}

The natural indices of the rational hierarchy are nonnegative integers
attached separately to the two target poles. Their finite algebra is
proved first. A different construction later extends the exponent of
the \emph{combined} quadratic denominator to one complex variable.
The distinction is essential: no recurrence with two independently
complex split indices is established here
\cite{VCDCMManuscript,VCDCMDiscussion08}.

\subsection{Rational traces and a legitimate infinite-limit passage}

On a counted pair chart use $w_j=\rho_j-c$, $Q_j=w_j^2-D$ for
each exterior occurrence. Strict exterior separation on compact
subcharts, together with \eqref{agc:pair-count}, makes
\begin{equation}
 S_n=\sum_jQ_j^{-n}\quad(n\ge1),\qquad
 T_n=\sum_jw_jQ_j^{-n}\quad(n\ge2),\qquad M=S_1
 \label{agc:hier-ST}
\end{equation}
absolute. Their respective large-$|\rho|$ tails at fixed orders are
$O(N^{1-2n}\log N)$ and $O(N^{2-2n}\log N)$.
The first odd order $T_1$ is deliberately absent.

For either weight $k(z)=Q_*^{-n}$ or
$k(z)=(z-c)Q_*^{-n}$ in these ranges,
\begin{equation}
 \sum_j k(\rho_j)=-\frac1{2\pi\ii}
                       \oint_\Gamma k(z)B(t,z)\,\dd z.
 \label{agc:hier-trace}
\end{equation}
Indeed, for one exterior pole $k(z)/(z-\rho)$ the sum of residues
at the target poles is $-k(\rho)$, since there is no residue at
infinity. Product regularizer constants multiply $\oint k=0$.
The target contribution $kQ_{*,z}/Q_*$ has all its poles enclosed
and decays faster than $z^{-1}$, so its combined residue also vanishes.
This proves \eqref{agc:hier-trace} by locally uniform product
convergence. The contour must enclose both targets; separating their
residues can introduce nonzero terms which cancel only in the sum.

Here is the precise passage used for the hierarchy. At a reference
slice, truncate the opposite-paired product with all conjugates,
ties and multiplicities retained, and include the target pair. Give
each resulting polynomial its own heat evolution initialized at that
slice. Its spatial jets converge on the cleared contour. Its first
time logarithmic jet is determined there by
\begin{equation}
 B_t=-\frac{H_{zzz}}H+\frac{H_zH_{zz}}{H^2}
     =-B_{zz}-2BB_z.
 \label{agc:hier-Burgers}
\end{equation}
Differentiating the fixed-contour coordinate and trace formulas,
including the moving kernels, therefore passes their first time jets
to the limit. At a finite multiple-root slice one first perturbs the
factors to simple roots and uses the same contour argument. Thus a
finite polynomial identity passes to the complete field once every
nonabsolute first-power sum has cancelled and only fixed-order
absolute moment products remain. This slice argument asserts no
common time interval for the polynomial approximants and no
uniformity as their moment order grows.

\subsection{The integer split recurrence}

Where $D\ne0$ choose a local analytic $s=\sqrt D$, and put
$h=2s$, $X_j=w_j-s$, $Y_j=w_j+s=X_j+h$. Define
\begin{equation}
 P_{p,q}=\sum_jX_j^{-p}Y_j^{-q},\qquad
 p,q\in\mathbb Z_{\ge0},\quad p+q\ge2.
 \label{agc:hier-P}
\end{equation}
These are absolute even on the axes. Their tails are
$O(N^{1-p-q}\log N)$ at fixed indices and fixed cleared targets.
For $m,n\ge1$ write
\[
 C_{m,n}=\sum_{i=2}^{m}P_{i,0}P_{m+1-i,n}
          +\sum_{j=1}^{n-1}P_{m,j}P_{0,n+1-j},
\]
and on an axis use
$C_{m,0}=\sum_{i=2}^{m-1}P_{i,0}P_{m+1-i,0}$ and
$C_{0,n}=\sum_{j=2}^{n-1}P_{0,j}P_{0,n+1-j}$.
An empty sum is zero. With $A_h=2/h-hM$, the recurrence is
\begin{align}
 P_{p,q}'={}&p\{C_{p+1,q}-(p+3)P_{p+2,q}
                  -(q+2)P_{p+1,q+1}+A_hP_{p+1,q}\}\notag\\
 &+q\{C_{p,q+1}-(p+2)P_{p+1,q+1}
                  -(q+3)P_{p,q+2}-A_hP_{p,q+1}\},
          &&p,q\ge1, \label{agc:hier-mixed}\\
 P_{p,0}'={}&p\{C_{p+1,0}-(p+3)P_{p+2,0}-2P_{p+1,1}
                                     +(2/h)P_{p+1,0}\},
          &&p\ge2, \label{agc:hier-axis-one}\\
 P_{0,q}'={}&q\{C_{0,q+1}-2P_{1,q+1}-(q+3)P_{0,q+2}
                                     -(2/h)P_{0,q+1}\},
          &&q\ge2. \label{agc:hier-axis-two}
\end{align}
Also $h'=4/h-2hM$. These identities hold through exterior
multiplicities on separated target charts.

\begin{proof}
Work first with finitely many simple roots. The target velocities are
$z_+'=2/h-2P_{1,0}$, $z_-'=-2/h-2P_{0,1}$, where the first-order
symbols are used only at this finite stage. Since
$X_j^{-1}-Y_j^{-1}=h/Q_j$, their difference gives $h'$.
An exterior root receives the target force
$2(X_j^{-1}+Y_j^{-1})$. For
$f_{m,n}(w)=(w-s)^{-m}(w+s)^{-n}$, the exterior--exterior force
contributes
\[
 U_{m,n}=-2\sum_{j<k}
             \frac{f_{m,n}(w_j)-f_{m,n}(w_k)}{w_j-w_k}
\]
to the block multiplied by $p$ or $q$ after differentiating its
weight. Reciprocal-power subtraction gives the finite identity
\begin{equation}
 -\frac{f_{m,n}(u)-f_{m,n}(v)}{u-v}
 =\sum_{i=1}^{m}f_{i,0}(u)f_{m+1-i,n}(v)
   +\sum_{j=1}^{n}f_{m,j}(u)f_{0,n+1-j}(v).
 \label{agc:hier-split-DD}
\end{equation}
To verify it, subtract the first factors while holding
$(v+s)^{-n}$ fixed, then subtract the second factors while
holding $(u-s)^{-m}$ fixed. The full left side is symmetric in
$u,v$, although this decomposition need not look symmetric.
Thus $U_{m,n}$ is the ordered off-diagonal sum on the right.
Its diagonal subtractions are $-mP_{m+1,n}-nP_{m,n+1}$.
For $m,n\ge1$, its two endpoint products are
$(P_{1,0}+P_{0,1})P_{m,n}$, and its remaining products are $C_{m,n}$.

The differentiated first block is
\[
 p\{(2/h-2P_{1,0})P_{p+1,q}
                  -2P_{p+2,q}-2P_{p+1,q+1}+U_{p+1,q}\};
\]
the second is the same expression with prefactor $q$, target
velocity $-2/h-2P_{0,1}$, and raised second index. Adding the
off-diagonal formulas cancels the common first-order sums; their
remaining difference is $P_{1,0}-P_{0,1}=hM$. The diagonals and
the two explicit target-force terms give exactly
\eqref{agc:hier-mixed}. On an axis both endpoint products contain
the same first-order sum and cancel it completely, leaving
\eqref{agc:hier-axis-one}--\eqref{agc:hier-axis-two}.
After these cancellations every surviving series is absolute.
The weight in \eqref{agc:hier-P} has the same fixed-contour trace
as \eqref{agc:hier-trace}. The preceding slice-jet passage therefore
proves the complete equations, including exterior multiplicities.
\end{proof}

Rows attached to different adjacent edges are constrained by their
shared divisor. If $P_{p,q}^{(n)}$ denotes \eqref{agc:hier-P} for
$(x_n,x_{n+1})$, restoring the omitted shared neighbor gives, for $r\ge2$,
\begin{equation}
 P_{r,0}^{(n)}+(-g_n)^{-r}
       =P_{0,r}^{(n+1)}+g_{n+1}^{-r}.
 \label{agc:hier-overlap}
\end{equation}
Both sides equal $\sum_{k\ne n+1}(x_k-x_{n+1})^{-r}$.
Within one row, $Y-X=h$ also gives
$hP_{p,q}=P_{p,q-1}-P_{p-1,q}$ for $p,q\ge1$,
$p+q\ge3$. At $(1,1)$ the difference must be combined before
summation. These necessary compatibility identities do not assert
that an arbitrary moment array comes from one divisor.

\subsection{Symmetric variables and their collision-regular evolution}

The split labels exchange when $s$ is replaced by $-s$. Their
regular symmetric form is \eqref{agc:hier-ST}, since
\begin{equation}
 S_n=P_{n,n},\qquad
 T_n=\tfrac12(P_{n-1,n}+P_{n,n-1}).
 \label{agc:hier-change}
\end{equation}
Conversely, put $r=\max(p,q)$ and write
\[
 X^{-p}Y^{-q}
     =\frac{(w-s)^{r-p}(w+s)^{r-q}}{(w^2-D)^r}.
\]
Its numerator has degree at most $2r-2$. Replacing even powers
$w^{2k}$ by $(Q+D)^k$ and odd powers by $w(Q+D)^k$ expresses
the result using only $S_n$ for $n\ge1$ and $T_n$ for $n\ge2$,
with polynomial coefficients in $s$. This proves equivalence of the
integer descriptions in their convergent ranges.

The collision-regular equations are
\begin{align}
 S_n'={}&2n\sum_{k=1}^{n}S_kS_{n+1-k}
       +2n\sum_{k=2}^{n}(T_kT_{n+2-k}+DS_kS_{n+2-k})\notag\\
       &-4n(n+2)S_{n+1}-4n(n+3)DS_{n+2},&&n\ge1,
                                                  \label{agc:hier-S}\\
 T_n'={}&2(2n-1)\sum_{k=2}^{n}T_kS_{n+1-k}
       +4nD\sum_{k=2}^{n}T_kS_{n+2-k}\notag\\
       &-4(n^2+n-1)T_{n+1}-4n(n+3)DT_{n+2},&&n\ge2.
                                                  \label{agc:hier-T}
\end{align}

For completeness, the symmetrization giving these coefficients can
be checked without manipulating a singular split label. At a finite
simple slice put $P=T_1=\sum_jw_j/Q_j$. The relative root equation is
\begin{equation}
 w_j'=4w_j/Q_j+2P+2\sum_{k\ne j}(w_j-w_k)^{-1}.
 \label{agc:hier-relative}
\end{equation}
For $A=(u^2-D)^{-1}$, $B_0=(v^2-D)^{-1}$, and
$L_r=\sum_{i=0}^{r-1}A^{r-1-i}B_0^i$ with $L_0=0$, elementary
subtraction gives
\begin{equation}
 \begin{split}
 \frac{A^r-B_0^r}{u-v}&=-(u+v)AB_0L_r,\\
 \frac{uA^r-vB_0^r}{u-v}
       &=-AB_0L_{r-1}-(D+uv)AB_0L_r.
 \end{split}
 \label{agc:hier-regular-DD}
\end{equation}
Here $AB_0L_r=\sum_{k=1}^r A^kB_0^{r+1-k}$. Summing over
unordered exterior pairs and subtracting the diagonals therefore gives
\begin{align*}
 -\sum_{i<j}\frac{w_iQ_i^{-n-1}-w_jQ_j^{-n-1}}{w_i-w_j}
 ={}&\frac12\sum_{k=1}^{n}(S_kS_{n+1-k}-S_{n+1})\\
 &+\frac D2\sum_{k=1}^{n+1}(S_kS_{n+2-k}-S_{n+2})\\
 &+\frac12\sum_{k=1}^{n+1}
                (T_kT_{n+2-k}-S_{n+1}-DS_{n+2}),\\
 -\sum_{i<j}\frac{Q_i^{-r}-Q_j^{-r}}{w_i-w_j}
 ={}&\sum_{k=1}^r(T_kS_{r+1-k}-T_{r+1}).
\end{align*}
In these two finite identities only, $T_1=P$ is allowed.
Differentiating $Q^{-n}$ using \eqref{agc:hier-relative} gives
\[
 nD'S_{n+1}-8n(S_{n+1}+DS_{n+2})-4nPT_{n+1}
       -4n\sum_{i<j}\frac{w_iQ_i^{-n-1}-w_jQ_j^{-n-1}}{w_i-w_j}.
\]
The $k=1,n+1$ odd endpoint products cancel $P$, and the
$D$-weighted endpoints cancel $-4nDMS_{n+1}$ from $nD'S_{n+1}$.
The diagonal coefficients left are $-4n(n+2)$ and $-4n(n+3)$,
proving \eqref{agc:hier-S}. For the odd weight use
\[
 (wQ^{-n})_w=(1-2n)Q^{-n}-2nDQ^{-n-1},\qquad
 (wQ^{-n})_D=nwQ^{-n-1}.
\]
The second divided-difference sum then cancels both terms involving
$P$; its $D$-weighted last endpoint cancels the interaction in
$nD'T_{n+1}$. The remaining direct and diagonal coefficients are
$-4(n^2+n-1)$ and $-4n(n+3)$, giving \eqref{agc:hier-T}.
The fixed-order slice passage proves both complete laws, through
target and exterior multiplicities.

In particular,
\begin{equation}
 \begin{split}
 M'&=2M^2-12S_2-16DS_3,\\
 S_2'&=8MS_2-32S_3+4DS_2^2-40DS_4+4T_2^2,\\
 T_2'&=(6M+8DS_2)T_2-20T_3-40DT_4.
 \end{split}
 \label{agc:hier-low}
\end{equation}
The spatial term $4T_2^2$ has a positive sign. All signs refer to
forward physical time; backward transport must reverse the time
derivative of the whole equation. These are componentwise identities
of the given complete field. They do not prove well-posedness for
arbitrary moment arrays or stability after a finite-order closure.

At a target merger the regular moments remain analytic, but a raw
split coordinate generally has the form $A(t)+s(t)C(t)$, with
$A,C$ analytic. Its coefficient at the double is
$C(t_*)=(p-q)\sum_jw_j(t_*)^{-p-q-1}$.
Since $D'(t_*)=2$, a nonzero coefficient produces a derivative of
order $1/s$. Thus collision regularity belongs to the symmetric
description, not to each separately labelled pole coordinate.

\subsection{The full inverse-power family and its Mellin meaning}

Choose a fixed length $\ell>0$. For each exterior occurrence select
a spatial logarithm of $Q(t,w)=w^2-D(t)$ near its location and set
\begin{equation}
 f_p(t,w)=\exp\{p(2\log\ell-\operatorname{Log}Q(t,w))\},\qquad
 Z(p,t)=\sum_j f_p(t,w_j),\qquad \Re p>\tfrac12.
 \label{agc:hier-Z}
\end{equation}
Then $Z(n,t)=\ell^{2n}S_n(t)$ for integer $n\ge1$.
The infinite inverse-power family retains all denominator scales;
increasing a real order emphasizes smaller exterior denominators.
For positive $Q_1<Q_2$, the ratio of their order-$n$ contributions
is $(Q_2/Q_1)^n$. This compares contributions within one order,
not dimensionful values of different moments. Nonreal denominators
carry phases and need not provide a positive moment sequence.

When every $Q_j>0$, define the denominator heat trace
$\Theta(v,t)=\sum_j\exp(-vQ_j/\ell^2)$ for $v>0$. The gamma
integral gives the genuine Mellin representation
\begin{equation}
 Z(p,t)=\frac1{\Gamma(p)}\int_0^\infty
                         v^{p-1}\Theta(v,t)\,\dd v,
                         \qquad \Re p>\tfrac12.
 \label{agc:hier-Mellin}
\end{equation}
Indeed, for $a=\Re p>1/2$, the integral of the sum of absolute
values is $\Gamma(a)\sum_j(\ell^2/Q_j)^a<\infty$.
Tonelli followed by Fubini proves the formula. It also holds with
the compatible logarithm if all $Q_j$ lie in a common closed sector
$\Re Q_j\ge\eta|Q_j|>0$, since the same argument uses
$(\Re Q_j)^{-a}$. A general complex exterior need not satisfy this
sector condition. In that case \eqref{agc:hier-Z} is the primary
complex-power trace; an unrotated positive-ray Laplace integral
must not be asserted for it.

\subsection{Local logarithms and the complete complex-time law}

For this analytic extension work at real times with a
conjugation-invariant pair, so $c,D$ are real. Assume, in addition to the complete
product, a fixed strip $|\Im\rho|\le A_0$, the projected unit count
$\#\{\rho:v\le\Re\rho\le v+1\}\le C\log(2+|v|)$,
and safe-line bounds
$|B(t,x\pm\ii Y)|\le C\log(2+|x|)$ with $Y>A_0$.
These are explicit source inputs, not consequences of a moment
recurrence. Keep $c,D$ in a compact counted chart and choose a large
centered cutoff $R_0$ so that $\Re(w^2-D)\ge(\Re w)^2/2$ for
$|\Re w|\ge R_0$, $|\Im w|\le Y$. Use the principal logarithm
on those two half-strips. Enclose the finitely many central exterior
clusters by disjoint cleared disks avoiding $Q=0$. On each disk
choose one analytic logarithm of $Q$ and continue it in time;
shrinking the chart makes this possible even through an exterior
merger. There is a chart constant $\Theta_0$ bounding the absolute
imaginary parts of these logarithms.

Let $\mathcal S$ denote $(2\pi\ii)^{-1}$ times integration over
the positive central disk boundaries and the boundaries of the two
infinite half-strips. On the latter the lower rays run rightward,
the upper rays leftward, the finite right connector downward and
the finite left connector upward. The central target pair is excluded.
For $h(t,w)=B(t,c(t)+w)$, the residue theorem gives
\begin{equation}
 \begin{split}
 Z&=\mathcal S(f_ph),\\
 Z_t&=\mathcal S((f_p)_t h+f_ph_t)\\
    &=\mathcal S\big(\{(f_p)_t-(f_p)_{ww}-c'(f_p)_w\}h
                                      +(f_p)_wh^2\big),\\
 h_t&=-h_{ww}-2hh_w+c'h_w,\qquad
 (f_p)_t=pD'f_p/Q.
 \end{split}
 \label{agc:hier-complex-time}
\end{equation}
The derivatives of $f_p,h$ here hold $w$ fixed.

We give the convergence details because a formal sum of zero
velocities at infinitely many multiple roots would not justify this
law. At fixed $1/2<b\le\Re p\le r$, counting and the chosen phases give
\begin{equation}
 |Z(p,t)|\le C_{b,r}\exp(\Theta_0|\Im p|),\qquad
 |Z_t(p,t)|\le C_{b,r}(1+|p|)^2\exp(\Theta_0|\Im p|)
 \label{agc:hier-p-envelope}
\end{equation}
on compact subcharts. For the value estimate, the tails are bounded
by $\sum_{n\ge1}n^{-2b}\log(2+n)$; central terms are finite and
separated. On the safe rays the absolute reciprocal-square and
reciprocal-cube sums yield $|h_w|+|h_{ww}|=O(\log(2+|x|))$.
Thus $h_t=O(\log^2(2+|x|))$ by \eqref{agc:hier-complex-time}.
Moreover
\[
 (f_p)_w=-2pw f_p/Q,\qquad
 (f_p)_{ww}=\{-2p/Q+4p(p+1)w^2/Q^2\}f_p.
\]
Both differentiated integrands consequently have an integrable
majorant $C(1+|p|)^2|x|^{-2b}\log^2(2+|x|)
\exp(\Theta_0|\Im p|)$. Cleared finite pieces have bounded
spatial jets and a positive $|Q|$ floor, giving the same exponent
envelope with a finite constant.

To remove the distant vertical closures, choose their real projections
in successive unit intervals at distance at least
$c/\log(2+|x|)$ from all zero projections. Such cuts exist by
deleting the short forbidden intervals and using the unit count.
Comparing $h$ on each cut with its value on the upper safe line
via the once-subtracted logarithmic derivative gives
$h=O(\log^2(2+|x|))$ there: the near unit shell pays inverse
projection clearance and the far shells have inverse-square decay.
The bounded-length connector integral is therefore
$O(|x|^{-2b}\log^2(2+|x|))$ and tends to zero. This proves the
spatial residue formula. The displayed majorants justify time
differentiation and normal convergence in $p$ for $\Re p>1/2$.
Finally, the difference between the two time integrands in
\eqref{agc:hier-complex-time} is the spatial derivative of
$-f_ph_w+(f_p)_wh-f_ph^2+c'f_ph$. Its circle integrals vanish,
its connector endpoint values cancel, and its values at infinite
ray endpoints tend to zero. Integration by parts proves the last
representation, including the full interaction term $(f_p)_wh^2$.

Two analytic time charts can differ on finitely many central clusters
by logarithm shifts $2\pi\ii k_C$. On their overlap,
\begin{equation}
 Z_B-Z_A=\sum_C(e^{-2\pi\ii k_Cp}-1)Z_{A,C}.
 \label{agc:hier-branch-transition}
\end{equation}
The same relation differentiates in time on the local cluster
contours. All corrections and their time derivatives vanish at
integer $p$. When real exterior denominators are positive, as for
a count-two disk containing the target segment, conjugate clusters
can be assigned conjugate logs, including opposite arguments
$+\pi,-\pi$ at nonreal conjugate locations with negative-real $Q$. This gives
$Z(p)=\overline{Z(\bar p)}$. If instead every negative-real
$Q_j$ is assigned $+\pi$ on a single slice, the exact defect is
\[
 Z_+(p)-\overline{Z_+(\bar p)}
   =-2\ii\ell^{2p}\sin(\pi p)\sum_{Q_j<0}|Q_j|^{-p}.
\]
The exceptional sum is finite. Resetting arguments at every time
need not define an analytic time chart. These branch restrictions
are essential away from the integer hierarchy.

\subsection{Damped recovery of the forcing}

The complex representation recovers the inherited forcing exactly.
Set $\delta=D/\ell^2$, $m=\ell^2M$, $\sigma=(t-t_0)/\ell^2$,
fix $\epsilon>0$, and use
\[
 K_\epsilon(p,t)=\frac{12e^{\epsilon(p-2)^2}}{p-2}
              +\frac{16\delta(t)e^{\epsilon(p-3)^2}}{p-3},
 \qquad G_\epsilon(p)=\frac{e^{\epsilon(p-1)^2}}{p-1}.
\]
For upward vertical lines $V_q$, define the two-line difference
$\mathcal V_{b,r}(F)=(2\pi\ii)^{-1}
(\int_{V_r}F\,\dd p-\int_{V_b}F\,\dd p)$. Then
\begin{equation}
 \begin{split}
 F(t)&=\mathcal V_{1,4}(K_\epsilon Z)
       =12Z(2,t)+16\delta Z(3,t)
       =\ell^4(12S_2+16DS_3),\\
 m(t)&=\mathcal V_{3/4,3/2}(G_\epsilon Z)=Z(1,t),\\
 \frac{\dd m}{\dd\sigma}&=2m^2-F,\qquad
 \frac{\dd\delta}{\dd\sigma}=2-4m\delta.
 \end{split}
 \label{agc:hier-forcing-contour}
\end{equation}
To prove this, close a counterclockwise rectangle between the chosen
verticals. The damping equals one at each residue and has modulus
$e^{\epsilon((\Re p-n)^2-(\Im p)^2)}$.
The horizontal integrals therefore vanish as their height tends
to infinity by \eqref{agc:hier-p-envelope}; the vertical integrals
are absolutely convergent. The same Gaussian dominates the time
derivative, including
$16(D'/\ell^2)e^{\epsilon(p-3)^2}Z/(p-3)$ from differentiating
the kernel. Residues give the first two formulas, and
\eqref{agc:pair-law}, \eqref{agc:hier-low} give the last two.
On a branch overlap the correction \eqref{agc:hier-branch-transition}
vanishes at the residue indices. The apparent poles of its product
with either kernel are removable, so the complete two-line values
and their derivatives agree across charts.

These contours reorganize the complete forcing. They provide neither
its sign nor an estimate stronger than evaluating the same integer
moments from the same information. In particular one cannot replace
the integer upper limits in \eqref{agc:hier-mixed} or
\eqref{agc:hier-S} by arbitrary complex indices. The single-complex-$p$
law \eqref{agc:hier-complex-time} is the proved analytic extension;
the independent canonical input needed to use the forcing for
continuation is considered next.

% 19 November 2026. Prepared for Edward Baker with Codex assistance.
% Lead gpt-6-astra/max. Existing comparisons are rederived, not new claims.
\section{What a backward estimate must control}
\label{sec:continuation}

\subsection{The cumulative drift budget}

The pair law of Section~\ref{sec:pair-chain} reads
\begin{equation}
 D_t=2-4DM,\qquad
 D_u=-\mathcal E,\qquad
 \mathcal E=2-4DM=2-g^2M.
 \label{cont:drift}
\end{equation}
The target is a time integral of the actual coupled quantity, not a
sign inferred from the approximation $d^2M\approx2$. Even
$g^2M=2+O(\varepsilon)$ does not select a sign for $D_u$; it selects
an error scale whose accumulated effect must be compared with the
initial squared gap.

\begin{proposition}[Cumulative continuation margin]
\label{cont:cumulative}
Suppose a counted pair satisfies \eqref{cont:drift} on a stopped path
$0\le u<\tau$, $D>0$, and $D(0)\ge d_0>0$. If an integrable real
function $\varepsilon$ satisfies $\mathcal E(u)\le\varepsilon(u)$,
then
\begin{equation}
 D(u)\ge d_0-\int_0^u\varepsilon(v)\dd v.
 \label{cont:integral}
\end{equation}
If the supplied inputs persist up to a possible first target merger
and
\begin{equation}
 \sup_{0\le s\le\Delta}\int_0^s\varepsilon(v)\dd v<d_0,
 \label{cont:strict}
\end{equation}
that merger cannot occur by $\Delta$. For a nonnegative allowance it
suffices to check the integral at $\Delta$.
\end{proposition}

\begin{proof}
Integrate \eqref{cont:drift}. The proposed first merger has $D\to0$,
whereas \eqref{cont:strict} leaves a uniform positive lower bound.
No differential equation is divided by $D$ at the event.
\end{proof}

This elementary comparison specifies a quantitative task. It does not
supply $\varepsilon$, a root-count tube, or coverage of other gaps.
When the allowance is signed, checking only its final cumulative value
could miss an earlier exhaustion; the supremum in
\eqref{cont:strict} is essential. At height $X$, with
$L=\log(X/(4\pi))$ and $D(0)$ of order $L^{-2}$, even a small
absolute drift error must be integrated on that scale.

\subsection{What initial interaction alone buys}

In an all-real simple configuration, exterior terms for an adjacent
pair are positive. Put
\begin{equation}
 \mathcal F=12S_2+16DS_3.
 \label{cont:F}
\end{equation}
The first moment equation gives
$M_u=-2M^2+\mathcal F\ge-2M^2$.
If $b_0=M(0)\ge0$, comparison with the scalar Riccati solution gives
$M(u)\ge b_0/(1+2b_0u)$. Retaining this lower bound in
$D_u=-2+4DM$ and using the factor $(1+2b_0u)^{-2}$ yields
\begin{equation}
 D(u)\ge(1+2b_0u)
       \{D(0)-2[1-b_0D(0)]u\}.
 \label{cont:initial-mass}
\end{equation}
The comparison is valid before its all-real stopped regime ends.
When $b_0D(0)<1$, its strict positive duration is
\begin{equation}
 u<\frac{D(0)}{2[1-b_0D(0)]}.
 \label{cont:initial-horizon}
\end{equation}
This derivation also covers $b_0=0$ by continuity. For an initial
near-clock configuration with $b_0D(0)\to1/2$, the duration approaches
$D(0)=g(0)^2/4$. It still vanishes at high height. This is the allowance
of this particular comparison, not a dynamical prediction of collision
\cite[\S\S13--15]{VCDCMManuscript}.

For example, an exact cosine heat field has stationary equally spaced
zeros, $M=2/g^2$ and $D=g^2/4$. It survives indefinitely, while
\eqref{cont:initial-horizon} stops at $g^2/4$. The forcing discarded
in the Riccati comparison balances $2M^2$ exactly. This comparison
control has uniform density; the canonical reference in
Section~\ref{sec:pair-chain} remains logarithmically height dependent.

\subsection{Keeping the joint forcing}

Define
\begin{equation}
 V=1-2DM,\qquad K=2D^2\mathcal F
                      =24D^2S_2+32D^3S_3.
 \label{cont:VK}
\end{equation}
These quantities combine gap size and exterior moments before an
inequality discards their dependence. Direct differentiation gives
\begin{equation}
 D_u=-2V,\qquad V_u=\frac{1-V^2-K}{D}.
 \label{cont:VK-law}
\end{equation}
Indeed $V_u=4M-4DM^2-2D\mathcal F$, and substitution of
$2DM=1-V$ gives the second equation. This is a regular calculation
for $D>0$; the comparison below is stopped before a merger.

\begin{theorem}[Joint forcing barrier]
\label{cont:barrier}
Suppose the real complete pair laws hold on a stopped path with $D>0$,
$V(0)\le v_0$, and $K(u)\ge1-\eta^2$ for a constant $\eta\ge0$.
For $v\ge\max\{v_0,\eta\}$,
\begin{equation}
 V(u)\le v,\qquad D(u)\ge D(0)-2vu.
 \label{cont:barrier-bound}
\end{equation}
A positive initial lower bound $d_0$ protects the target pair through
$\Delta$ if $2v\Delta<d_0$ and the stipulated inputs persist up to
a first event.
\end{theorem}

\begin{proof}
On each compact pre-event interval the coefficient below is continuous.
For $w=V-v$, \eqref{cont:VK-law} implies
\[
 w_u+\frac{V+v}{D}w
 \le\frac{\eta^2-v^2}{D}\le0.
\]
A positive integrating factor and $w(0)\le0$ prove $w\le0$.
Integrate $D_u=-2V$, then exhaust the stopped path by compact
subintervals. A first merger contradicts the resulting strictly positive
lower bound as in Proposition~\ref{cont:cumulative}.
\end{proof}

The proof itself needs no positive summands: a signed complete bound
for $K$ in a conjugation-invariant counted tube also suffices. Positive
omission of a tail, however, requires the all-real adjacent-pair regime.
If another pair collides first and ends that regime, it ends that use of
positivity. Protection of one pair is not protection of the complete
population.

For $D(0)\asymp L^{-2}$, sufficient rates to pay a fixed duration are
\begin{equation}
 V(0)_+=o(L^{-2}),\qquad
 \sup_u(1-K(u))_+=o(L^{-4}).
 \label{cont:rates}
\end{equation}
The square root of the forcing deficiency explains the stronger second
rate. These are sufficient requirements of this barrier, not necessary
conditions for actual noncollision.

\subsection{Finite heads: an exact comparison and its perturbation}

For an equally spaced comparison configuration with target endpoints
$0,g$, the $m$th exterior occurrence on either side has
$Q_m=m(m+1)g^2$. With $D=g^2/4$ and $J$ such occurrences on each side,
let $M_J,K_J$ denote their contributions. Telescoping gives
\begin{equation}
 2DM_J=\frac{J}{J+1},\qquad
 K_J=1-\frac1{(J+1)^3}.
 \label{cont:lattice-head}
\end{equation}
For the second identity the contribution from both sides is
\[
 \frac3{m^2(m+1)^2}+\frac1{m^3(m+1)^3}
   =\frac1{m^3}-\frac1{(m+1)^3}.
\]
These exact formulas are local comparison tools, not a uniform-density
model of the complete zeta divisor.

They have a useful robust version. In the all-real adjacent-pair stopped
regime, suppose the actual target gap obeys
$g\ge(1-e)d$, $0\le e<1/3$, and the $m$th selected exterior zero
on each side is at distance $r$ from the nearer endpoint, with
\begin{equation}
 r\le(1+2e)md,\qquad r+g\le(1+2e)(m+1)d,
 \quad 1\le m\le J.
 \label{cont:head-input}
\end{equation}
Writing $\gamma=(1-e)/(1+2e)$ gives
$D/[r(r+g)]\ge\gamma^2/[4m(m+1)]$.
Since the quadratic and cubic powers in $K$ are nonnegative,
\begin{equation}
 K\ge1-18e-(J+1)^{-3},\qquad
 V\le6e+(J+1)^{-1}.
 \label{cont:head-bounds}
\end{equation}
To check the constants, use $1-\gamma\le3e$,
$1-\gamma^6\le18e$ and $1-\gamma^2\le6e$ in
\eqref{cont:lattice-head}, and omit only positive exterior terms.
A uniform version of the first bound and the initial instance of the
second supply Theorem~\ref{cont:barrier}. Adding more initial neighbors
alone does not provide the whole-time first bound.

\subsection{What the existing canonical acquisition already uses}

The source analysis in \cite[\S15]{VCDCMManuscript}, based on
\cite[Theorem 1.5 and \S9]{Polymath2019}, works on a fixed slab
$0<t_0\le t\le t_*\le1/2$. The height onset depends exponentially
on $1/t_0$. The positive-head acquisition is restricted to the globally
all-real simple stopped portion, begun from an all-real simple configuration. On that same slab it first locates individual consecutive
zeros near the source clock; it then estimates complete $M$ and a
neighboring head for $K$. With $L=\log(X/(4\pi))$, a head of
order $L^3$ neighbors and its tail payment give
\[
 |V|=O(L^{-3}),\qquad (1-K)_+=O(L^{-9})+o(L^{-m})
\]
for each fixed $m$, after the sufficiently large height is chosen.
The error terms inherited from localization are smaller than every
fixed inverse power of $L$ on this fixed positive slab. They do not
have that property uniformly as its lower time tends to zero.

For completeness, the mechanism is the preceding head calculation.
On a spatial radius $L^2$ the clock geometry has normalized error
$O(LX^{-c t_0}+L^4/X)$; hence \eqref{cont:head-bounds} with
$J\asymp L^3$ gives the forcing floor. The complete inverse-square
exterior contributes $O(L^{-3})$ to the normalized interaction.
Inserting these quantities into \eqref{cont:barrier-bound} gives a
relative loss $O(\Delta/L)$ for fixed $\Delta$ within the supplied
slab. The localization theorem already guarantees reality and simplicity
in that exterior. Consequently this acquisition explains the dynamics
and the available rates, but does not extend the zero-free/collision-free
coverage beyond the imported theorem. The new question is precisely
how to acquire useful inputs after that localization estimate ceases
to provide them.

% 19 November 2026. Drafted for Edward Baker with OpenAI Codex assistance.
% Helper explicitly dispatched as gpt-6.1-sol, reasoning effort ultra;
% independent serving-backend attestation was not available to the helper.
\section{The full arithmetic reference and its useful domain}
\label{sec:arithmetic-source}

The arithmetic input to a backward continuation need not be reduced to
the two leading carriers that produce an asymptotic clock.  At a positive
time the finite Riemann--Siegel approximation retains many arithmetic
terms, with their phases, before its error is estimated.  The question
for this program is whether those terms can supply effective gap data
where the two-carrier argument no longer does so.  This is a question
about an additional quantitative estimate.  The approximation and the
identities below do not themselves extend the known all-real region.
We use the published approximation of \cite{Polymath2019} and the
selected-source analysis in \cite{VCDCMDiscussion11,VCDCMDiscussion12}.

\subsection{A holomorphic finite reference}

To avoid confusing an arithmetic carrier with the complete exterior
interaction of a zero pair, write $\mathfrak M_t$ for the function called
$M_t$ in \cite{Polymath2019}.  On
$\mathcal D=\C\setminus(-\infty,1]$, use the principal logarithms to set
\begin{align}
 L_0(s)&=\operatorname{Log}s+\operatorname{Log}(s-1)
       -\frac{s}{2}\log\pi+\log\frac{\sqrt{2\pi}}{16}
       +\left(\frac{s}{2}-\frac12\right)
           \operatorname{Log}\frac{s}{2}-\frac{s}{2},
       \label{as:carrier-log}\\
 \mathfrak M_0(s)&=\exp L_0(s),\qquad
 \alpha(s)=L_0'(s)
       =\frac{1}{2s}+\frac{1}{s-1}
         +\frac12\operatorname{Log}\frac{s}{2\pi},
       \label{as:alpha}\\
 \mathfrak M_t(s)&=\mathfrak M_0(s)
                  \exp\left(\frac{t}{4}\alpha(s)^2\right).
       \label{as:carrier}
\end{align}
These carriers are nonvanishing on $\mathcal D$.  The logarithm in
\eqref{as:carrier-log} fixes an actual analytic function, rather than a
pointwise choice of argument.  For $n\in\N$ define
\begin{align}
 \ell_n(s)&=L_0(s)-s\log n,\qquad
 a_n(s)=\ell_n'(s)=\alpha(s)-\log n,\\
 p_n(t,s)&=\exp\left(\ell_n(s)+\frac{t}{4}a_n(s)^2\right).
       \label{as:pn}
\end{align}
Expanding the square gives the useful exact identity
\begin{equation}
 p_n(t,s)=\mathfrak M_t(s)\,
     \exp\left(\frac{t}{4}\log^2 n\right)
     n^{-s-(t/2)\alpha(s)}.
 \label{as:pn-dirichlet}
\end{equation}
In particular, the $n=1$ term is $\mathfrak M_t(s)$, and all the terms
with $2\le n\le N$ remain part of the reference.

Put
\begin{equation}
 s_+(z)=\frac{1+\ii z}{2},\qquad
 s_-(z)=\frac{1-\ii z}{2},\qquad
 P_N(t,z)=\sum_{n=1}^N\{p_n(t,s_+(z))+p_n(t,s_-(z))\}.
 \label{as:PN}
\end{equation}
For a fixed integer $N$, this is holomorphic in $z$ on the right
half-plane.  Indeed $\Re z>0$ places the two $s$-arguments in opposite
open half-planes, away from the cut of $L_0$.  For real $t$ it satisfies
$P_N(t,\overline z)=\overline{P_N(t,z)}$ there.  It therefore has the
local conjugation symmetry required of a reference for a real zero pair.
This assertion about its definition is separate from the domain of a
published error bound.

For $z=x+\ii y$ with $x\ge200$ and $0\le y\le1$, the convenient
nonvanishing normalizer is
\begin{equation}
 \mathcal B_t(z)=\mathfrak M_t(s_-(z)),\qquad
 \gamma_t(z)=\frac{\mathfrak M_t(s_+(z))}
                  {\mathfrak M_t(s_-(z))},\qquad
 b_n^t=\exp\left(\frac{t}{4}\log^2 n\right).
 \label{as:normalizer}
\end{equation}
Consequently the normalized full reference is
\begin{equation}
 f_{t,N}(z)=\frac{P_N(t,z)}{\mathcal B_t(z)}
   =\sum_{n=1}^N\frac{b_n^t}{n^{s_-(z)+(t/2)\alpha(s_-(z))}}
    +\gamma_t(z)\sum_{n=1}^N
       \frac{b_n^t}{n^{s_+(z)+(t/2)\alpha(s_+(z))}}.
 \label{as:full-f}
\end{equation}
This form makes holomorphy visible.  A formula written using $x,y$ and
complex conjugates must first be rewritten in this form before complex
spatial differentiation is justified; see \cite[\S8, (92)]{Polymath2019}.

\subsection{The exact arithmetic terms behind the approximation}

The finite reference comes from an exact expansion, rather than an
independent model of the zero divisor.  Define
\begin{equation}
 K(s)=\frac{s(s-1)}{16}\pi^{-s/2}\Gamma(s/2),\qquad
 r_{0,n}(s)=K(s)n^{-s}.
 \label{as:exact-initial-term}
\end{equation}
For nonreal $s$ and $t>0$ the Gaussian evolution of this term is
\begin{equation}
 r_n(t,s)=\frac1{\sqrt\pi}\int_{\R}
          r_{0,n}(s+\sqrt t\,v)e^{-v^2}\,\dd v.
 \label{as:exact-term}
\end{equation}
The horizontal integration line avoids the poles of the meromorphic
factor $K$.  On compact sets separated from the real $s$-axis, its
growth along that line is slower than a Gaussian, so the integral and
the differentiations used below converge locally uniformly.

For completeness, the exact Riemann--Siegel remainder before Gaussian
evolution is
\begin{equation}
 R_{0,N}(s)=K(s)\int_{N\swarrow N+1}
       \frac{w^{-s}e^{\pi\ii w^2}}
            {e^{\pi\ii w}-e^{-\pi\ii w}}\,\dd w.
 \label{as:exact-RS-remainder}
\end{equation}
Here the oriented diagonal line has direction $e^{5\pi\ii/4}$ and
crosses the real axis between $N$ and $N+1$; the conventional logarithm
along the Riemann--Siegel contour is understood.  Let $R_N(t,s)$ be its
Gaussian evolution as in \eqref{as:exact-term}, and write
$F^*(s)=\overline{F(\overline s)}$.  The exact decomposition is
\begin{align}
 H_t(z)={}&\sum_{n=1}^N r_n(t,s_+(z))
        +\sum_{n=1}^N r_n^*(t,s_-(z))\notag\\
        &+R_N(t,s_+(z))+R_N^*(t,s_-(z)).
 \label{as:exact-finite-expansion}
\end{align}
It holds for every fixed nonnegative integer $N$, $t>0$, and
$\Re z\ne0$; this is the Gaussian form of the Riemann--Siegel identity
in \cite[\S4]{Polymath2019}.  The star is retained to display the
reflection in the exact formula.  The coefficients and the Gaussian
kernel are real under conjugation, so that $r_n^*=r_n$ on the two
nonreal components.  Replacing each exact $r_n$ by \eqref{as:pn} gives
\eqref{as:PN}; neither replacement removes the complete remainder.

The approximation of an individual term has an explicit, relative
source bound.  With $s=\sigma+\ii T$, $T>10$, and $0<t\le1/2$, one may
write
\begin{equation}
 r_n(t,s)=p_n(t,s)(1+\delta_n(t,s)),\qquad
 |\delta_n(t,s)|\le
 \exp\left(
 \frac{(t^2/8)|\alpha(s)-\log n|^2+t/4+1/6}{T-3.33}
 \right)-1.
 \label{as:term-source-bound}
\end{equation}
The corresponding lower-half-plane statement follows by conjugation.
This is the term estimate of \cite[Proposition 6.1]{Polymath2019}; its
error is relative to an individual carrier term, rather than to $H_t$
near one of its zeros.  Summing it and paying the Riemann--Siegel
remainder supplies the complete approximation error.

The standard cutoff in that complete bound is
\begin{equation}
 N(t,x)=\left\lfloor\sqrt{\frac{x}{4\pi}+\frac{t}{16}}\right\rfloor.
 \label{as:natural-cutoff}
\end{equation}
Theorem 1.3 of \cite{Polymath2019} gives, throughout
\begin{equation}
 0<t\le\tfrac12,\qquad x\ge200,\qquad 0\le y\le1,
 \label{as:source-domain}
\end{equation}
the complete bound
\begin{equation}
 H_t(z)=P_{N(t,x)}(t,z)+E_0(t,z),\qquad
 |E_0(t,z)|\le |\mathcal B_t(z)|
             \{e_A(t,z)+e_B(t,z)+e_{C,0}(t,z)\}.
 \label{as:source-value-bound}
\end{equation}
The three $e$'s in this display denote the explicit positive budgets of
that theorem, with its definitions (71)--(74); they are a named imported
input, not quantities bounded in this manuscript.  The first two pay
the term replacements and the last pays the whole remainder when its
leading correction is omitted.  No numerical conclusion below requires
an unevaluated choice of these budgets.  An application must use their
actual formulas, with the same normalization and domain.  Bounds below
the real axis can be supplied by conjugating the physical approximation.

Two qualifications matter when using \eqref{as:source-value-bound} in
a differential argument.  First, the theorem does not impose
$x\ge\exp(C/t)$.  Its domain is larger than the clock regime.  Second,
the natural cutoff jumps.  One must use a time and spatial cell on
which \eqref{as:natural-cutoff} has one value, or keep a fixed $N$ and
explicitly transfer the omitted or added terms to the error.  For example,
at the same $(t,z)$,
\begin{equation}
 P_{N+1}(t,z)-P_N(t,z)
       =p_{N+1}(t,s_+(z))+p_{N+1}(t,s_-(z)).
 \label{as:cutoff-jump}
\end{equation}
This is a known change of reference.  It must be included in a seam
payment if a cell crosses a cutoff surface.  Differentiating the floor
in \eqref{as:natural-cutoff} as if it were constant would instead lose
that payment.

\subsection{What the two-carrier clock discards}

On the real axis put $L=\log(x/(4\pi))$ and $m=\log n$.
Directly from \eqref{as:alpha},
$\Re\alpha((1-\ii x)/2)=L/2+O(x^{-2})$.  Therefore the magnitude of
an individual normalized term in the first sum of \eqref{as:full-f} is
\begin{equation}
 \left|b_n^t n^{-s_--(t/2)\alpha(s_-)}\right|
   =n^{-1/2}
      \exp\left\{-\frac{t}{4}m(L-m)+O(tm/x^2)\right\}.
 \label{as:clock-term-decay}
\end{equation}
For $n\le N(t,x)$, one has
$m\le L/2+O(t/x)$.  When $tL$ is sufficiently large, summing
\eqref{as:clock-term-decay} over $n\ge2$ gives a negative power of
$x$ with exponent proportional to $t$, after decreasing the absolute
constant in that exponent.  The same is true for the reflected sum.
The essential large parameter for this dominance is $t\log x$.
Accordingly a uniform small-tail argument has an onset of the form
$x\ge\exp(C/t)$, with an absolute constant that has to be large enough
for all the other estimates as well.

The coefficient $t/4$ in \eqref{as:clock-term-decay} follows from the
definitions by expanding a square.  It is not replaced here by the
$t/2$ coefficient appearing in an intermediate display of the source's
\S9; the discrepancy and its source check are recorded in
\cite{VCDCMDiscussion11}.  We extract no numerical onset constant from
this qualitative summation argument.

In the dominance region the reference reduces to
\begin{equation}
 P_N(t,z)\simeq\mathfrak M_t(s_+(z))+
                      \mathfrak M_t(s_-(z)).
 \label{as:two-carriers}
\end{equation}
Cancellation of these two carriers yields the clock
\begin{equation}
 G(x,t)=\frac{x}{4\pi}\log\frac{x}{4\pi}
       -\frac{x}{4\pi}+\frac{11}{8}
       +\frac{t}{16}\log\frac{x}{4\pi},\qquad
 G(q_n(t),t)=n.
 \label{as:clock}
\end{equation}
Theorem 1.5 of \cite{Polymath2019} supplies actual real simple zeros
near these clock points above its height and index onsets.  The existing
analysis \cite{VCDCMManuscript} uses one buffered onset on a slab
$0<t_0\le t\le t_*\le1/2$, and hence one set of common labels.
The distinction between a height onset and an index onset is preserved;
both must be met before those labels can be used.

The full reference \eqref{as:PN} remains available when the $n\ge2$
terms cannot be treated as a vanishing tail.  They can then change the
phase, location, and separation of its zeros.  They belong in the
reference field being followed.  Treating them as a new small error
without a bound would simply reintroduce the clock hypothesis.  On the
other hand, retaining them does not prove that their own divisor is
real and simple, or that the source error is smaller than a gap margin.
The useful initial datum is an enclosure of an \emph{actual} gap
$g=b-a$, or of $D=g^2/4$, with isolation and completeness paid.  A mean
spacing $4\pi/L$ becomes such a datum only after localization errors
are smaller than the relevant separation.  This is why a source formula
with a wider domain is an opportunity for continuation, rather than an
already acquired continuation theorem.

\subsection{The strip certificate and the existing upper bounds}

A global upper bound for $\Lambda$ has another consumer besides a list
of separated real pairs.  We recall the exact domains of the criterion
being compared.  In the nontrivial regime $0<t_0<1/2$, $X>0$ and
$0<y_0<\sqrt{1-2t_0}$, the three inputs of
\cite[Theorem 1.2]{Polymath2019} are the following zero exclusions:
\begin{align}
 &\zeta(\sigma+\ii T)\ne0
 &&\left(\frac{1+y_0}{2}\le\sigma\le1,
                       \quad0\le T\le X/2\right),
                       \label{as:strip-initial}\\
 &H_{t_0}(x+\ii y)\ne0
 &&\left(x\ge X+\sqrt{1-y_0^2},
                   \quad y_0\le y\le\sqrt{1-2t_0}\right),
                       \label{as:strip-final}\\
 &H_t(x+\ii y)\ne0
 &&\left\{\begin{aligned}
       &X\le x\le X+\sqrt{1-y_0^2},\quad0\le t\le t_0,\\
       &\sqrt{y_0^2+2(t_0-t)}\le y\le\sqrt{1-2t}.
       \end{aligned}\right.
                       \label{as:strip-barrier}
\end{align}
Together they give
\begin{equation}
 \Lambda\le t_0+\frac{y_0^2}{2}.
 \label{as:strip-output}
\end{equation}
These parameter restrictions keep the displayed regions meaningful
and include the applications discussed below. The last two inputs
are nonvanishing problems for the complete source, so finite arithmetic
lower bounds and their full error budgets can be useful without resolving
every real gap.

The extra $y_0^2/2$ is the payment for contracting a global zero strip
to the real axis.  In particular, \eqref{as:strip-final} and
\eqref{as:strip-barrier} do not assert that all zeros are already real
at $t_0$.  A real-pair backward argument beginning at $t_0$ cannot use
that strip statement as an all-real, separated initial divisor.  A later
all-real slice and the quantitative separation needed by its chosen
comparison are additional data.  Even a global all-real statement by
itself is not a numerical lower bound for every gap.

Platt and Trudgian prove finite RH verification through
$T_{\rm PT}=3\,000\,175\,332\,800$.  Their Corollary 2 obtains
$\Lambda\le0.2$ by applying the existing Polymath parameter table to
that verified height \cite{PlattTrudgian2021}.  The conversion is $X=2T$.
The result supplies complete finite zero exclusion at time zero, rather
than the fine individual gap enclosures proposed here.

At the pinned revision of Gomila's proof package, the reported row is
\begin{equation}
 X=6\,000\,000\,185\,827,\qquad t_0=\frac{129}{800},\qquad
 y_0^2=\frac{87677}{2500000}.
 \label{as:gomila-row}
\end{equation}
Its claimed computer-assisted unconditional result, awaiting peer
review, is $\Lambda\le0.1787854$ \cite{Gomila2026Proof}.  The exact
parameter arithmetic is
\begin{equation}
 t_0+\frac{y_0^2}{2}=0.16125+0.0175354=0.1787854,
 \qquad
 T_{\rm PT}-\frac X2=\frac{350479773}{2}>0.
 \label{as:gomila-arithmetic}
\end{equation}
Thus the exact verified endpoint covers the requested initial height.
Using the rounded headline $3\cdot10^{12}$ would obscure that reserve.
This check is parameter bookkeeping, not an independent replay of the
proof package.  Its finite arithmetic lower bounds, infinite-tail
argument, and barrier over the intervening time interval are separate
certificate obligations.  The repository audit has a stated partial
replay scope \cite{NewmanCertificateAudit}; neither that audit nor the
present exposition upgrades the external package's review status.

The live explanatory page \cite{Gomila2026Blog} also reports a later
0.1782354 calculation at the same verified height, with its mathematical
review still pending. That separate research report is not a bound
verified by this investigation and is not substituted for the pinned
package compared here.

These comparisons specify what a future gain must accomplish.  A local
improvement in a gap continuation margin can be useful, but it is not
yet an improvement of \eqref{as:strip-output}.  The latter also needs
complete height coverage and the correct initial, final, and barrier
interfaces.  Conversely a better strip certificate need not supply the
separated all-real starting data of a gap program.  The two objectives
share arithmetic source information and pay different final margins.

% 19 November 2026. Drafted for Edward Baker with OpenAI Codex assistance.
% Helper explicitly dispatched as gpt-6.1-sol, reasoning effort ultra;
% independent serving-backend attestation was not available to the helper.
\section{The forced residual and the joint gap readout}
\label{sec:residual-transport}

We next distinguish an exact heat evolution from its explicit arithmetic
approximation.  This distinction determines which error may be
propagated.  All calculations in this section use a fixed cutoff, or
one cell on which the cutoff is constant.  They are elementary
consequences of the source representation and the physical equation
$H_t=-H_{zz}$.  They supply an interface for a quantitative estimate;
no such new canonical estimate is asserted here.

\subsection{The exact term evolves by heat; the explicit term has a defect}

For the exact Gaussian term \eqref{as:exact-term}, differentiation and
integration by parts give
\begin{align}
 \partial_t r_n(t,s)
  &=\frac1{2\sqrt{\pi t}}\int_{\R}
       v\,r_{0,n}'(s+\sqrt t\,v)e^{-v^2}\,\dd v\notag\\
  &=\frac1{4\sqrt\pi}\int_{\R}
       r_{0,n}''(s+\sqrt t\,v)e^{-v^2}\,\dd v
    =\frac14\partial_s^2r_n(t,s).
 \label{rt:exact-term-heat}
\end{align}
The sub-Gaussian growth described after \eqref{as:exact-term} kills
the boundary term and justifies the differentiations on compact
nonreal $s$-sets and positive-time intervals.  The same calculation
applies to the exact remainder where its Gaussian formula is used.

There is no corresponding exact heat assertion for $p_n$.  Suppress
the index $n$, and write
\begin{equation}
 a=\alpha-\log n,\qquad b=\alpha',\qquad c=\alpha'',\qquad
 p=\exp(\ell_n+ta^2/4).
 \label{rt:abc}
\end{equation}
These are functions of $s$; explicitly
\[
 b(s)=-\frac1{2s^2}-\frac1{(s-1)^2}+\frac1{2s},\qquad
 c(s)=\frac1{s^3}+\frac2{(s-1)^3}-\frac1{2s^2}.
\]
Differentiating this actual exponential, rather than a remainder
symbol, gives
\begin{align}
 \frac{p_t}{p}&=\frac{a^2}{4},\qquad
 \frac{p_s}{p}=a+\frac t2ab,\notag\\
 \frac{p_{ss}}p&=\left(a+\frac t2ab\right)^2
                   +b+\frac t2(b^2+ac).
 \label{rt:pn-jets}
\end{align}
Subtracting one fourth of the last expression from the first gives
the exact defect
\begin{equation}
 \mathcal R_n(t,s):=\left(\partial_t-\frac14\partial_s^2\right)p_n(t,s)
 =-\frac{p_n(t,s)}4
   \left\{b+ta^2b+\frac t2(b^2+ac)
                       +\frac{t^2}{4}a^2b^2\right\}.
 \label{rt:term-defect}
\end{equation}
For the additive term error $d_n=r_n-p_n$ this implies
\begin{equation}
 \left(\partial_t-\frac14\partial_s^2\right)d_n=-\mathcal R_n.
 \label{rt:term-error-equation}
\end{equation}
The term $b$ in \eqref{rt:term-defect} is present even at $t=0$:
the explicit expression has replaced both the gamma factor and its
Gaussian evolution by an approximation.  A small value of the defect
is not a bound for the solution of \eqref{rt:term-error-equation}
until its data and domain are specified.

Because $\partial_z^2p_n(t,s_\pm(z))=-p_{n,ss}(t,s_\pm(z))/4$,
the physical fixed-$N$ reference has defect
\begin{equation}
 \mathcal R_N(t,z):=\partial_tP_N(t,z)+\partial_z^2P_N(t,z)
   =\sum_{n=1}^N\{\mathcal R_n(t,s_+(z))
                      +\mathcal R_n(t,s_-(z))\}.
 \label{rt:physical-defect}
\end{equation}
Define the actual error by $E=H-P_N$.  Then, exactly,
\begin{equation}
 E_t+E_{zz}=-\mathcal R_N.
 \label{rt:physical-error}
\end{equation}
For backward physical time $u=t_*-t$, with every quantity evaluated
at $t_*-u$, this becomes the forward forced heat equation
\begin{equation}
 E_u=E_{zz}+\mathcal R_N.
 \label{rt:backward-error}
\end{equation}
The sign and coefficient here follow from both changes of variables:
$s_\pm=(1\pm\ii z)/2$ and $u=t_*-t$.  There is no missing factor of
four in the physical equation.

Using the sum of the exact $r_n$ terms as the reference would instead
give zero defect, by \eqref{rt:exact-term-heat}; its full remainder
would still need to be controlled.  Using the explicit $p_n$ terms
gives a directly differentiable formula and the known forcing
\eqref{rt:physical-defect}.  These are two different acquisition
interfaces, not a reason to discard the remainder in either one.

\subsection{What is required to propagate the residual}

Equation \eqref{rt:backward-error} on a bounded spatial neighborhood
does not determine an effective error enclosure there from a local
initial bound alone.  For example, for $A>0$ the entire functions
\[
 E_K(u,z)=\varepsilon\exp\{K(z-A)+K^2u\},\qquad K>0,
\]
satisfy $E_u=E_{zz}$ and
$\sup_{-A\le x\le A}|E_K(0,x)|\le\varepsilon$, whereas
$E_K(u,A)=\varepsilon e^{K^2u}$.  Thus the same local initial upper
bound permits arbitrarily large later values.  The uncontrolled
growth outside the interval is exactly what prevents a uniform
conclusion.  This example concerns the information in a local heat
bound, not the canonical source.

One possible propagation statement is the following deliberately
conditional version.  Its strong global hypotheses are displayed
because they have not been acquired for the arithmetic reference.

\begin{proposition}[A whole-strip residual bound]
\label{rt:duhamel}
Let $0\le u\le U$ and $h>0$.  Suppose that $E$ and $R$ are
holomorphic in $z$ on a neighborhood of the closed horizontal strip
$|\Im z|\le h$, continuous in time, and that $E$ is a classical
solution of $E_u=E_{zz}+R$ there.  Suppose that $E$ is uniformly
bounded on this whole time-strip, that
\begin{equation}
 \sup_{x\in\R,\,|\eta|\le h}|E(0,x+\ii\eta)|\le\varepsilon_0,
 \qquad
 \sup_{x\in\R,\,|\eta|\le h}|R(v,x+\ii\eta)|\le r(v),
 \label{rt:global-data}
\end{equation}
where $r$ is integrable.  Assume the usual continuity needed to take
these bounded classical solutions to their initial data.  Then
\begin{equation}
 \sup_{x\in\R,\,|\eta|\le h}|E(u,x+\ii\eta)|
       \le\varepsilon(u):=\varepsilon_0+\int_0^u r(v)\,\dd v.
 \label{rt:duhamel-bound}
\end{equation}
For $0<h'<h$, Cauchy's estimate also gives
\begin{equation}
 \sup_{x\in\R,\,|\eta|\le h'}|E_{zz}(u,x+\ii\eta)|
       \le\frac{2\varepsilon(u)}{(h-h')^2}.
 \label{rt:duhamel-jets}
\end{equation}
\end{proposition}

\begin{proof}
For each fixed $\eta$, put
$G_u(q)=(4\pi u)^{-1/2}\exp(-q^2/(4u))$.  The bounded whole-line
heat solution has the Duhamel representation
\begin{align*}
 E(u,x+\ii\eta)={}&\int_{\R}G_u(x-\xi)E(0,\xi+\ii\eta)\,\dd\xi\\
 &+\int_0^u\int_{\R}G_{u-v}(x-\xi)
                      R(v,\xi+\ii\eta)\,\dd\xi\,\dd v.
\end{align*}
The heat kernel has integral one and is nonnegative on the real line.
Taking absolute values proves \eqref{rt:duhamel-bound}, uniformly in
$\eta$.  The representation follows by the fundamental solution and
uniqueness in the bounded classical class; it can also be verified by
regularizing in time and passing to the bounded initial data.  Apply
Cauchy's formula on disks of radius smaller than $h-h'$ and then take
their radius to $h-h'$ to obtain \eqref{rt:duhamel-jets}.
\end{proof}

The hypotheses include every real $x$, not just the part of the strip
where a zero pair is enclosed.  In particular this proposition is not
automatically applicable to the branch-defined reference $P_N$ of
\eqref{as:PN}, which was defined on a right half-plane.  A genuine use
must provide an admissible extension and the whole-line bounds, or use
a bounded-domain heat problem with certified lateral boundary data.
If a Gaussian integral is truncated, its tails must be paid as well.
For example, a global bound $M$ on its initial integrand pays at most
$M\operatorname{erfc}(A/(2\sqrt u))$ outside $|x-\xi|\le A$;
without such a bound, the displayed tail allowance has no premise.

Normalization does not remove this obligation.  If $H=\mathcal B F$
with a nonvanishing analytic factor, then $F$ satisfies
\begin{equation}
 F_t+F_{zz}+2\frac{\mathcal B_z}{\mathcal B}F_z
       +\frac{\mathcal B_t+\mathcal B_{zz}}{\mathcal B}F=0.
 \label{rt:normalized-equation}
\end{equation}
Thus a bound for the normalized error $H/\mathcal B-f$ cannot be
propagated by the plain unit-coefficient heat kernel while ignoring
the two additional coefficients.  Moreover a small absolute physical
error can be large relative to $H$ near a zero.  The contour denominator
needed next is a separate quantitative datum.

\subsection{A two-zero contour retains a joint error}

Let $\Gamma$ be a positively oriented piecewise smooth simple closed
contour, fixed over the time cell.  Assume that $H$ and the reference
$P$ are holomorphic on a neighborhood of its closed interior, nonzero
on a neighborhood of $\Gamma$, and each has exactly two zeros inside,
counted with multiplicity.  In this subsection
$P$ can be $P_N$ or another explicitly specified analytic reference.
The exact heat law is imposed on $H$, while
\begin{equation}
 R=P_t+P_{zz},\qquad E=H-P
 \label{rt:general-reference}
\end{equation}
is retained for $P$.  Counts and contour clearance must hold throughout
the cell; they are not inferred from a sampled value.

For $F=H$ or $P$, define
\begin{equation}
 A_j^F=\frac1{2\pi\ii}\oint_\Gamma z^j\frac{F_z}{F}\,\dd z,
 \qquad c_F=\frac{A_1^F}{2},\qquad
 D_F=\frac{A_2^F}{2}-\frac{(A_1^F)^2}{4}.
 \label{rt:pair-moments}
\end{equation}
If the two zeros are $a,b$, the argument principle gives
$c_F=(a+b)/2$ and $D_F=(b-a)^2/4$.  These contour expressions remain
defined when the two occurrences merge.  For a conjugation-invariant
pair, $D_F$ is real; it is positive for distinct real zeros and
negative for a distinct nonreal conjugate pair.

Differentiate the logarithmic derivative of $H$ and integrate by parts:
\begin{equation}
 (A_j^H)_t=\frac{j}{2\pi\ii}\oint_\Gamma
                          z^{j-1}\frac{H_{zz}}H\,\dd z,
 \qquad
 (D_H)_t=\frac1{2\pi\ii}\oint_\Gamma
                          (z-c_H)\frac{H_{zz}}H\,\dd z.
 \label{rt:exact-gap-readout}
\end{equation}
For example,
$\partial_t(H_z/H)=\partial_z(H_t/H)
=-\partial_z(H_{zz}/H)$, which proves the first formula.  Substitution
in the definition of $D_H$ proves the second.  This is an exact readout
of the pair law, including at a merger, rather than a new exclusion
criterion.

On the cleared contour, elementary algebra gives
\begin{equation}
 \frac{H_{zz}}H-\frac{P_{zz}}P
       =\frac{P E_{zz}-P_{zz}E}{P(P+E)}.
 \label{rt:joint-quotient}
\end{equation}
The two error terms in the numerator share the same residual.  For a
spatially constant multiplier $E=\lambda(t)P$, the numerator vanishes
exactly.  Separate bounds for $P E_{zz}$ and $P_{zz}E$ lose that
dependence.  This identity identifies a possible place to retain
arithmetic cancellation, but it proves no bound for that cancellation
in the canonical source.

There are two further payments.  The actual center in
\eqref{rt:exact-gap-readout} is $c_H$, not $c_P$.  Also $P_{zz}/P$
alone does not give the time motion of the zeros of $P$.  Repeating
the preceding calculation with $P_t=-P_{zz}+R$ gives
\begin{equation}
 (D_P)_t=\frac1{2\pi\ii}\oint_\Gamma
                   (z-c_P)\left(\frac{P_{zz}}P-\frac RP\right)\,\dd z.
 \label{rt:reference-gap-readout}
\end{equation}
Subtracting it from \eqref{rt:exact-gap-readout} yields the full identity
\begin{align}
 (D_H-D_P)_t=\frac1{2\pi\ii}\oint_\Gamma
 \bigg\{&(z-c_H)\frac{P E_{zz}-P_{zz}E}{P(P+E)}
       +(c_P-c_H)\frac{P_{zz}}P\notag\\
       &+(z-c_P)\frac RP\bigg\}\,\dd z.
 \label{rt:full-gap-error}
\end{align}
In particular, paying only \eqref{rt:joint-quotient} and calling its
integral the error in $(D_P)_t$ would omit both a center error and the
known heat defect.

The center can itself be enclosed without differentiating a value
error.  Suppose $|E/P|\le q<1$ on a neighborhood of $\Gamma$.
The convergent branch $\log(1+E/P)$ is then single-valued there.
Integration by parts along the contour gives
\begin{equation}
 c_H-c_P=-\frac1{4\pi\ii}\oint_\Gamma
                          \log(1+E/P)\,\dd z,
 \qquad
 |c_H-c_P|\le\frac{|\Gamma|}{4\pi}\{-\log(1-q)\}.
 \label{rt:center-bound}
\end{equation}
Here $|\Gamma|$ is the contour length.  No logarithm is asserted through
the interior zeros of $H$ or $P$; only the collar of the contour is used.
The same smallness condition, together with the two-zero count for $P$,
also supplies the count for $H$ by Rouch\'e's theorem.

For a concrete bound on \eqref{rt:full-gap-error}, set on each slice
\begin{gather*}
 m=\inf_\Gamma|P|,\quad \varepsilon=\sup_\Gamma|E|<m,\quad
 J=\sup_\Gamma|P E_{zz}-P_{zz}E|,\\
 Q=\sup_\Gamma|P_{zz}/P|,\quad W=\sup_\Gamma|R/P|,\quad
 B_H=\sup_\Gamma|z-c_H|,\quad B_P=\sup_\Gamma|z-c_P|.
\end{gather*}
Then
\begin{equation}
 |(D_H-D_P)_t|\le\mathcal A(t):=
 \frac{|\Gamma|}{2\pi}
 \left\{\frac{B_HJ}{m(m-\varepsilon)}
              +|c_H-c_P|Q+B_PW\right\}.
 \label{rt:gap-error-bound}
\end{equation}
The $J$ bound may be obtained directly from a coupled enclosure, or
estimated less sharply from paid separate jets.  The latter is a valid
fallback, with no promised gain from correlation.  Every denominator
floor is displayed.  Integrating in backward time gives the conditional
comparison
\begin{equation}
 |D_H(u)-D_P(u)|\le |D_H(0)-D_P(0)|
                      +\int_0^u\mathcal A(t_*-v)\,\dd v.
 \label{rt:integrated-gap-error}
\end{equation}
To use this for continuation, the right side must be smaller than an
available positive lower bound for $D_P(u)$, with the counts, reality
symmetry, and all-time contour clearance retained.  Alternatively one
can bound the exact drift in \eqref{rt:exact-gap-readout} directly and
feed its integrated shrinking allowance to Section~\ref{sec:continuation}.
Neither use is supplied by a small residual norm in isolation.

\subsection{The center of a corrected approximation must also change}

The leading Riemann--Siegel correction gives two different value
approximations.  Schematically, with the source normalizer,
\begin{equation}
 \frac{H_t}{\mathcal B_t}=f_{t,N}+e_0,
 \qquad
 \frac{H_t}{\mathcal B_t}
       =\left(f_{t,N}-\frac{C_t}{\mathcal B_t}\right)+e_1.
 \label{rt:corrected-centers}
\end{equation}
The published bound for $e_1$ replaces the omitted-correction budget
$e_{C,0}$ by the smaller retained-correction budget $e_C$
\cite[\S6, (69)--(70)]{Polymath2019}.  A lower bound $|f_{t,N}|\ge L$
therefore does not automatically consume the $e_1$ allowance.  It needs
a lower bound for $|f_{t,N}-C_t/\mathcal B_t|$.  The unconditional
reverse-triangle transfer pays
\begin{equation}
 \left|f_{t,N}-\frac{C_t}{\mathcal B_t}\right|
          \ge L-\left|\frac{C_t}{\mathcal B_t}\right|,
 \label{rt:correction-payment}
\end{equation}
which can spend the margin recovered in the remainder.  This was the
applicability issue for the triangle estimate examined in
\cite{NewmanCertificateAudit}.  It is scoped to that estimate and its
information: it does not exclude phase-sensitive lower bounds or a
different arithmetic reorganization.

There is an additional analytic issue for a contour consumer.  A
displayed correction written in terms of $x,y$, a real part, and a
floor-valued cutoff is a value formula.  Before taking its complex
derivatives, one must establish a holomorphic representative on the
chosen neighborhood or use a different analytic reference with a
separately paid correction.  Simply assigning a smaller value error
to $P_N-C_t$ does not establish that reference's analyticity or heat
defect.  If such a reference is supplied, all of
\eqref{rt:physical-defect}--\eqref{rt:gap-error-bound} must use that
same reference, its center, and its actual defect.

Arithmetic mollifiers illustrate another legitimate cancellation that
requires a transfer.  For a finite Dirichlet polynomial
$A(s)=\sum_{n\le N}a_n n^{-s}$ and
$m(s)=\sum_{d\le D}\lambda_d d^{-s}$, exact multiplication gives
\begin{equation}
 m(s)A(s)=\sum_{r\le DN}
           \left(\sum_{dn=r}\lambda_d a_n\right)r^{-s}.
 \label{rt:mollifier-convolution}
\end{equation}
One can combine the coefficients in the inner sum before taking
absolute values.  If a full calculation establishes $|m f|\ge Q>0$
and $|m|\le M$, then $m\ne0$ and
\begin{equation}
 |f|\ge Q/M.
 \label{rt:mollifier-transfer}
\end{equation}
Gomila's finite-sum binding applies this type of arithmetic
reorganization to the two pieces of the normalized approximation and
returns a bound for the original $f_t$
\cite{Gomila2026Binding,VCDCMDiscussion12}.  The convolution and division
are mathematical steps with their own positive-margin hypotheses.

For a moving gap, multiplication is not a free change of the physical
equation.  If $P$ has defect $R$ and $m=m(t,z)$ is a nonvanishing
analytic factor on the relevant region, then
\begin{equation}
 (mP)_t+(mP)_{zz}
      =mR+(m_t+m_{zz})P+2m_zP_z.
 \label{rt:factor-defect}
\end{equation}
The factor preserves the enclosed zero divisor when it has no zeros
there, but it changes the heat defect and the second-logarithmic
quotient.  The zero dynamics are invariant only when these extra terms
are retained in the full readout.  Thus the static success of a
mollifier lower bound does not imply that the multiplied field solves
the same Newman heat equation.

Finally the generic residual interface must preserve actual collisions.
For the exact control $H=P=z^2-2(t-T)$ one has $E=R=0$, while the two
real zeros merge at $t=T$.  Here $D=2(t-T)$, so an integrated backward
shrinking allowance reaches the initial $D$ exactly at the merger;
the strict continuation margin fails.  Perfect approximation has not
protected a colliding reference.  The more detailed premise screen in
Section~\ref{sec:program} distinguishes the complete canonical
theta source from the colliding source controls.  No generic error
identity above supplies that distinction.  The open arithmetic task is
to acquire a complete, source-specific bound at the final gap or strip
consumer, and to compare it with direct evaluation using the same
information and error payments.

% 19 November 2026. Prepared for Edward Baker with OpenAI Codex assistance.
% Lead recorded gpt-6-astra/max. Program statements are prospective.
\section{Synthesis, controls and outstanding challenges}
\label{sec:program}

\subsection{The analytical information being carried forward}

The pair coordinates compress a local divisor without discarding its
complete exterior. The gap chain makes this exterior geometric: nearby
gaps strongly influence the response, while distant zeros enter through
summable dipole tails. The integer hierarchy keeps the shared divisor
constraints; the complex Mellin extension supplies another representation
with its own convergence and branch requirements. These are compatible
views of one physical field. An arbitrary vector of moments need not
come from that field, and evolving a truncated head without its boundary
rows need not preserve that compatibility.

The earlier manuscript proves a common-time approximation by initializing
one finite paired product and evolving it with the polynomial heat
operator \cite[\S\S6--8]{VCDCMManuscript}. Its fixed-input convergence,
counted-contour transfer and conditioning qualifications are relevant
background. The present program uses the arithmetic approximation as a
more structured candidate reference, rather than assuming that generic
product convergence already supplies a useful relative error at large
height. Both require a count, a denominator floor and a complete path.

The normalization issue has two different manifestations. If
$H=AQ$ locally with $A$ analytic and nonzero, a contour integral of
$\partial_z\log A$ contributes no divisor count. But
\[
 \frac{H_{zz}}{H_z}(a)
   =\frac{Q_{zz}}{Q_z}(a)+2\frac{A_z}{A}(a)
\]
at a simple zero $a$ of $Q$. Thus the same factor affects zero velocity
under the heat law. A spatially constant multiplier cancels from that
ratio; an evolving reference with such a multiplier still has to pay
its time defect. These observations explain the earlier localization
analogy without constructing a localization or index-theoretic proof.
A count of two is the same before and at a double zero, so count
invariance alone cannot protect a gap
\cite{VCDCMDiscussion07,VCDCMDiscussion08}.

\subsection{A common contour makes the comparison explicit}

A short derivation specifies how a field estimate becomes a gap estimate.
Let a positively oriented contour $\Gamma$ contain exactly two zeros,
with multiplicity, of each of two holomorphic fields $H$ and $P$, and
no zeros on $\Gamma$. Write their centers and discriminants as
$c_H,D_H$ and $c_P,D_P$. Suppose $H/P$ has a single-valued analytic
logarithm $h$ on a neighborhood of the contour, for example because
$|H-P|<|P|$ throughout that neighborhood. For the power sums
\[
 A_j(H)=\frac1{2\pi\ii}\oint_\Gamma z^j\frac{H_z}{H}\dd z,
 \quad c_H=A_1(H)/2,\quad D_H=A_2(H)/2-c_H^2,
\]
integration by parts gives
\begin{equation}
 A_j(H)-A_j(P)=-\frac{j}{2\pi\ii}
                 \oint_\Gamma z^{j-1}h(z)\dd z.
 \label{prog:log-moments}
\end{equation}
If
$I_0=(2\pi\ii)^{-1}\oint h\dd z$ and
$I_c=(2\pi\ii)^{-1}\oint(z-c_P)h\dd z$, then
\begin{equation}
 c_H-c_P=-\tfrac12 I_0,\qquad
 D_H-D_P=-I_c-\tfrac14 I_0^2.
 \label{prog:centered-error}
\end{equation}
To verify the second equality, substitute
$A_2(H)-A_2(P)=-2(2\pi\ii)^{-1}\oint zh\dd z$ and expand
$c_H^2=(c_P-I_0/2)^2$. The shared logarithm and full square remain.
Separately maximizing their contributions may give a valid but weaker
outer bound. A constant $h$ gives $I_0=I_c=0$ exactly.

Together with the joint quotient identity of
Section~\ref{sec:residual-transport}, this gives a concrete paired
comparison: direct contour evaluation and time-integrated drift must
use the same source information and maintain the same correlations.
These identities rederive the relevant transfer principle in
\cite[\S\S7,19--20]{VCDCMManuscript}. They are not evidence that an
iteration of the same constraints adds information.

\subsection{What a first research result should accomplish}

The first comparison should fix one reference field, one parameter domain
and one target functional before numerical refinement. For a real pair,
that functional may be the cumulative drift in
Proposition~\ref{cont:cumulative}. An alternative is the whole-time
forcing floor in Theorem~\ref{cont:barrier}. Their acquisition requirements
are different, and an estimate for one must not be silently substituted
for the other.

There are three connected analytical possibilities. First, estimate the
joint residual numerator directly so that common multiplicative errors
cancel before an absolute value is taken. Second, use the forced heat
representation or its adjoint to weight only those residual components
which reach the selected contour functional, while paying the full
spatial complement and boundary data. Third, retain the phases of the
arithmetic correction and its reference sum so that the corrected center
is bounded as one expression. These are proposed mechanisms, not acquired
canonical bounds. A mollifier that improves a value bound is not
automatically a stable reference for a derivative or time-dependent
quantity.

A successful finite result would be a smaller complete error allowance
or a longer certified interval for the fixed consumer, compared with
a direct method using the same information. A failure could also be
useful if it identifies the particular discarded correlation or growing
constant that prevents improvement. Such a failure is scoped to that
comparison. It cannot justify a universal claim that dynamics, different
mollifiers or other arithmetic decompositions are incapable of progress.

\subsection{Collision controls and failed premises}

The collision controls in \cite{RepositoryControls} distinguish a valid
transfer from an unjustified source-specific conclusion. The following
is an analytical premise screen, not a numerical replay of the control
certificates.

The simplest exact test is
\begin{equation}
 H(t,z)=P(t,z)=z^2-2(t-T),\qquad T>0.
 \label{prog:quadratic}
\end{equation}
It has zero approximation error and zero heat defect, but its real pair
merges at $t=T$. Here $M=K=0$, $V=1$ and
$D(u)=2(t_*-T)-2u$. The cumulative allowance is exhausted exactly at the
merger; the strict margin fails. Thus a sound residual theorem preserves
the event. Perfect approximation of a colliding reference is not
protection against collision.

An initial almost regular chain also need not provide a whole-time
forcing floor. For $k>0$ and $0<\alpha_0<1$, the exact backward heat
field
\begin{equation}
 F(u,z)=e^{-k^2u}\cos(kz)+\alpha_0
 \label{prog:cosine-control}
\end{equation}
has $F_u=F_{zz}$. Put $\alpha(u)=\alpha_0e^{k^2u}$ and
$\theta(u)=\arccos\alpha(u)$. Its short pair about $\pi/k$ has
\[
 D(u)=\theta(u)^2/k^2,\qquad
 V(u)=\frac{\theta(u)\alpha(u)}{\sqrt{1-\alpha(u)^2}},\qquad
 u_c=k^{-2}\log(1/\alpha_0).
\]
The formulas follow by solving $\cos(kz)=-\alpha(u)$ and differentiating.
As $\alpha_0\to0$, the initial gaps approach $\pi/k$ and
$V(0)\to0$, yet a merger occurs at $u_c$. For fixed $k$ the exterior
of the selected cluster stays separated there; its inverse powers
remain summable, so $K\to0$ as $D\to0$. The whole-time near-unit
forcing premise fails. This is a finite positive atomic source, not
theta, and serves only as a control on inference from initial geometry.

The remaining source controls fail different canonical premises:
\begin{itemize}
\item Euler-factor translates and weighted-square, finite-source,
finite-score and finite-prime perturbations alter the full normalized
theta arithmetic. Preserved finite measurements do not restore its
complete source identity. Documented endpoint-defective subclasses also
lack the imported strip.
\item Reflected super-Gaussian sources fail the order-one product input
used here as well as literal theta identification. Smooth shape and
positive-primitive controls lack that complete arithmetic identification;
positivity alone supplies no missing gap margin.
\item Atomic and finite-profile constructions fail the full theta profile.
Their event does not refute a theorem that uses the actual profile, but
it must survive any generic approximation theorem whose hypotheses they
satisfy.
\item RH-conditional endpoint constructions change ordinary Dirichlet
arithmetic and assume RH in their construction. Their escaping event
heights do not provide a fixed-compact obstruction.
\item Changed observation kernels fail the physical cosine/Newman
readout; fractional profiles fail integer divisor multiplicities; static
divisor surgery supplies no common source heat path. Strong-collar and
annular polynomial controls still test local algebra and the full square.
\end{itemize}
Identifying a missing canonical premise is not proof that a proposed
coarse estimate uses that premise effectively. Every future quantitative
claim must identify where the unchanged arithmetic actually enters.

\subsection{Outstanding analytical challenges}

The remaining work has a definite dependency order.

\begin{enumerate}
\item \emph{Acquire a useful reference and initial gap.}
The full arithmetic approximation is available beyond first-term clock
dominance, but its small remainder does not locate its own zeros.
A counted contour, initial separation and a compatible correction center
must be established without assuming the desired continuation.

\item \emph{Control the residual in the norm the gap uses.}
The forced error equation needs time-uniform spatial derivatives and
complete tail or boundary data. Its inverse denominator and center
errors must be paid. The main question is whether the common numerator,
its weighted time integral, or arithmetic phase information produces a
useful improvement after all of those costs.

\item \emph{Keep the changing environment coupled.}
The Mellin hierarchy organizes higher moments but is not a finite closure.
A source-constrained realization or direct source estimate must retain
shared zeros, moving gaps and complete exterior forces. At a regular
merger $K=2D^2\mathcal F$ vanishes, so assuming a whole-path positive
floor without acquiring it would already assume event exclusion.

\item \emph{Extend the source domain, not just its precision.}
At fixed positive time the favorable moment estimates inherit individual
clock separation. A new result must state exactly which time/height
domain it adds or which complete comparison it improves. The onset
$x\ge\exp(C/t)$ is a sufficient source threshold; its divergence at zero
is not removed by substituting an unevaluated constant or taking a
nonuniform limit.

\item \emph{Assemble local information into a global statement.}
A protected pair or finite block must be joined to compact and high-height
regions with no omitted gaps, interfaces or exits. If the Polymath strip
criterion is the eventual consumer, its initial, exterior and
intermediate-time barrier hypotheses and its $y_0^2/2$ time payment
remain. A local real-gap estimate cannot replace those obligations
without a new proved bridge.
\end{enumerate}

The purpose of the new track is to resolve these analytical dependencies
in sustained research sequences: derive a candidate estimate, test its
premises, compare its complete consequence, and use the result to choose
the next step. Numerical work supports those decisions through fixed
algebraic checks or explicitly certified examples. The unresolved
mathematical task is a source-specific estimate whose useful margin
persists beyond the information already supplied by direct localization.

% 19 November 2026. Drafted for Edward Baker with OpenAI Codex assistance.
% Lead actual backend/effort unavailable; helpers explicitly Sol6.1/ultra.
\section{Finite-cell propagation and a sharper complete allowance}
\label{sec:local-estimates}

The first continuation fixes the physical reference $P_N$ of
\eqref{as:PN} and compares two ways to acquire its residual on one
counted contour. The signed endpoint identity \eqref{prog:centered-error}
is inherited. The new estimates below fill in the lateral boundary
payment and sharpen the complete value-to-discriminant allowance. They
are finite estimates, not an RH-equivalent reformulation. Their arithmetic
application must still acquire the count, clearance and positive margin.

\subsection{One fixed arithmetic cell}

Let $0<U<t_*\le1/2$, $0<h\le1$, $X-A\ge200$, and require
\begin{equation}
 N^2\le\frac{X-A}{4\pi}+\frac{t_*-U}{16},\qquad
 \frac{X+A}{4\pi}+\frac{t_*}{16}<(N+1)^2.
 \label{le:cutoff}
\end{equation}
The natural cutoff is then the same integer throughout
$0\le u\le U$, $|\Re z-X|\le A$, $|\Im z|\le h$, with
$t=t_*-u$. This avoids a seam payment without discarding arithmetic terms.
It permits an interval of centers precisely when
$2A<4\pi(2N+1)-\pi U/4$, before imposing $X-A\ge200$.
Use the fixed circle $\Gamma:|z-X|=r<\min(A,h)$.
Neither its two-zero count nor a real positive reference discriminant is
assumed to follow from the cutoff inequalities.

For clarity, the complete published value budget can be written without
the normalizer. Put $T=x/2$, $\eta=|y|$, $\sigma_\pm=(1\pm\eta)/2$,
$w_\pm=\sigma_\pm+\ii T$, $T'=T+\pi t/8$, and
$a_{\rm RS}=\sqrt{T'/(2\pi)}$. Define
\begin{align}
 \epsilon_n(t,w)&=\exp\left(
 \frac{(t^2/8)|\alpha(w)-\log n|^2+t/4+1/6}{T-3.33}
 \right)-1,\notag\\
 \delta_\sigma(t,T)&=
 \left(\frac{0.397\,9^\sigma}{a_{\rm RS}-0.865}
                 +\frac5{3(T-6)}\right)
                 \exp\left(\frac{3.49}{T-4}\right).
 \label{le:source-errors}
\end{align}
Then Corollary~6.5 of \cite{Polymath2019} gives
\begin{align}
 |H_t(z)-P_N(t,z)|\le\mathcal V_N(t,x,\eta)
 :={}&\sum_{n=1}^N\sum_{\nu\in\{-,+\}}
                 |p_n(t,w_\nu)|\epsilon_n(t,w_\nu)\notag\\
 &+e^{t\pi^2/64}|\mathfrak M_0(\ii T')|
                   (1+\delta_{\sigma_-}+\delta_{\sigma_+}).
 \label{le:source-value}
\end{align}
The $1$ in the second line pays the leading omitted correction. It is
not replaced by the smaller retained-correction error. The source domain
is $0<t\le1/2$, $x\ge200$, $|y|\le1$ and $N=\lfloor a_{\rm RS}\rfloor$;
\eqref{le:cutoff} enforces this throughout the cell. The source's branch
names differ from ours; the explicit upper-half-plane arguments $w_\pm$
and conjugation fix that convention. The formula supplies initial and
lateral data as well as a direct endpoint envelope. The full finite
forcing remains \eqref{rt:physical-defect}, with its complex sum retained
before taking a modulus whenever it is evaluated jointly.

\subsection{Bounded-domain heat propagation with both boundaries}

Write $x_-=X-A$, $x_+=X+A$, $I=(x_-,x_+)$ and $L=2A$.
For every fixed $y\in[-h,h]$, $e(u,x)=E(u,x+\ii y)$ satisfies
$e_u=e_{xx}+\mathcal R_N$. Suppose the data and forcing are continuous,
and $E$ is $C^1$ in time, $C^2$ in $x$ and holomorphic on a neighborhood
of the closed rectangle, locally uniformly in time. These hypotheses
hold for the fixed finite reference and canonical field on this cell.
There is no whole-line extension premise.

Let $K(\tau;x,\xi)$ be the Dirichlet heat kernel on $I$:
\begin{equation}
 K(\tau;x,\xi)=\frac2L\sum_{n\ge1}e^{-(n\pi/L)^2\tau}
  \sin\frac{n\pi(x-x_-)}L\sin\frac{n\pi(\xi-x_-)}L,
 \qquad \tau>0.
 \label{le:kernel}
\end{equation}
It is nonnegative. Put
$k_-=\partial_\xi K|_{\xi=x_-}$ and
$k_+=-\partial_\xi K|_{\xi=x_+}$, both nonnegative. With
$f_y(\xi)=E(0,\xi+\ii y)$ and $g_\pm(v,y)=E(v,x_\pm+\ii y)$,
the exact formula is
\begin{align}
 E(u,x+\ii y)={}&\int_I K(u;x,\xi)f_y(\xi)\dd\xi\notag\\
 &+\int_0^u\int_I K(u-v;x,\xi)
             \mathcal R_N(t_*-v,\xi+\ii y)\dd\xi\dd v\notag\\
 &+\int_0^u\{k_-(u-v;x)g_-(v,y)
                    +k_+(u-v;x)g_+(v,y)\}\dd v.
 \label{le:bounded-duhamel}
\end{align}
Indeed, differentiate $\int_I K(u-v;x,\xi)E(v,\xi+\ii y)\dd\xi$.
Green's identity yields the boundary term
$[K E_\xi-K_\xi E]_{x_-}^{x_+}$ and the forcing integral. Integration
in $v$, followed by the approximate-identity limit at $u$, proves
\eqref{le:bounded-duhamel}. The variable $y$ is a parameter, so no top
or bottom boundary integral enters this one-dimensional heat problem.

Set $S(\tau,x)=\int_I K(\tau;x,\xi)\dd\xi$ and
$Q_\pm(\tau,x)=\int_0^\tau k_\pm(s;x)\dd s$. Applying the formula to
the constant solution gives
\begin{equation}
 S+Q_-+Q_+=1.
 \label{le:mass}
\end{equation}
The maximum principle, comparing the boundary responses with the
half-line heat solution, gives
\begin{equation}
 1-S(\tau,x)\le\min\left\{1,
 \operatorname{erfc}\frac{x-x_-}{2\sqrt\tau}
 +\operatorname{erfc}\frac{x_+-x}{2\sqrt\tau}\right\}.
 \label{le:boundary}
\end{equation}
The sum is an even convex function of $x-X$ in $|x-X|<A$.
Thus on the circle it is at most
\begin{equation}
 \beta_\Gamma(\tau)=\min\left\{1,
 \operatorname{erfc}\frac{A-r}{2\sqrt\tau}
 +\operatorname{erfc}\frac{A+r}{2\sqrt\tau}\right\},
 \qquad \beta_\Gamma(0)=0.
 \label{le:beta}
\end{equation}

If $|f_y|\le\varepsilon_0$, $|g_\pm|\le M_b$ and
$|\mathcal R_N(t_*-v,\xi+\ii y)|\le\rho(v)$ over their full domains,
positivity gives the paid pointwise bound
\begin{equation}
 |E(u,x+\ii y)|\le\varepsilon_0S(u,x)+M_b(1-S(u,x))
                +\int_0^u\rho(v)S(u-v,x)\dd v.
 \label{le:mass-bound}
\end{equation}
In particular, when $M_b\ge\varepsilon_0$,
\begin{equation}
 \sup_\Gamma|E(u,z)|\le\varepsilon_0+
          (M_b-\varepsilon_0)\beta_\Gamma(u)+\int_0^u\rho(v)\dd v.
 \label{le:scalar-heat}
\end{equation}
For $M_b<\varepsilon_0$ keep \eqref{le:mass-bound}, or discard the
negative boundary correction. Inserting an upper bound for $1-S$ into
a negative multiple would be invalid. Spatially varying data may instead
be integrated in \eqref{le:bounded-duhamel}; the scalar suprema are only
one relaxation. The equally informed direct calculation receives this
propagated envelope too, and can intersect it with \eqref{le:source-value}.

\subsection{Recenter the logarithm before bounding the divisor}

The next elementary estimate concerns the full value-error disk, not
the quadrature aliases treated in the parent manuscript's common-log
Parseval estimate \cite[\S22]{VCDCMManuscript}. It uses the same principle
of retaining shared Fourier information.

\begin{lemma}[A centered logarithm disk]\label{le:log-disk}
For $0\le q<1$ and $|w|\le q$, define
\[
 \lambda=-\log(1-q),\qquad
 \kappa=\tfrac12\log(1-q^2),\qquad
 \eta=\operatorname{arctanh}q.
\]
The principal logarithm obeys
$|\log(1+w)-\kappa|\le\eta$. For $q>0$, $\eta<\lambda$.
\end{lemma}
\begin{proof}
The case $q=0$ is immediate. Put $v=(1+w)/\sqrt{1-q^2}$.
Its disk has center $\cosh\eta$ and radius $\sinh\eta$.
The Cayley map $\xi=(v-1)/(v+1)$ sends this disk onto
$|\xi|\le\tanh(\eta/2)$. The identity
$\log v=2\operatorname{arctanh}\xi$, on the branch containing $v=1$,
and the absolutely convergent odd power series give
$|\log v|\le2\operatorname{arctanh}|\xi|\le\eta$.
The disk lies in the right half-plane, so this is precisely the
principal logarithm in the statement. Finally,
$\lambda-\eta=-\tfrac12\log(1-q^2)>0$.
\end{proof}

\begin{theorem}[A complete coupled circle allowance]\label{le:gap-bound}
Let $H,P$ be holomorphic on a neighborhood of $|z-X|\le r$, have
exactly two zeros there counted with multiplicity, and be nonzero on
the circle. Suppose $|H/P-1|\le q<1$ on it. Define $c_P,D_P,c_H,D_H$
by the contour moments, and put $d=|c_P-X|/r$. Then
\begin{equation}
 |D_H-D_P|\le r^2 F(\eta,d),\qquad
 F(\eta,d)=\max_{0\le x\le\eta}
       \left\{\sqrt{\eta^2-x^2}+dx+\frac{x^2}{4}\right\}.
 \label{le:F}
\end{equation}
Here $F(0,d)=0$. Useful explicit upper bounds are
\begin{align}
 F(\eta,d)&\le\eta\sqrt{1+d^2}+\eta^2/4,
                                                        \label{le:cauchy}\\
 F(\eta,d)&\le\eta+\frac{\eta d^2}{2-\eta}\quad(0<\eta<2).
                                                        \label{le:parabolic}
\end{align}
For the centered subcase,
\begin{equation}
 F(\eta,0)=
 \begin{cases}\eta,&0\le\eta\le2,\\1+\eta^2/4,&\eta\ge2.
 \end{cases}
 \label{le:centered}
\end{equation}
The maximum $F$ optimizes the displayed two-coefficient relaxation;
the Cauchy and parabolic bounds are further upper relaxations. No
sharpness over physical heat fields is asserted.
\end{theorem}
\begin{proof}
Let $\ell=\log(H/P)$ on the circle collar and write
$\ell(X+re^{\ii\theta})-\kappa=\sum_{j\in\mathbb Z}a_j e^{\ii j\theta}$.
The strict relative clearance supplies the collar branch. Bessel's
inequality and Lemma~\ref{le:log-disk} give
$|a_{-1}|^2+|a_{-2}|^2\le\eta^2$. Constants have zero contour moments.
Substituting the two Fourier coefficients into
\eqref{prog:centered-error} gives exactly
\begin{equation}
 D_H-D_P=-r^2a_{-2}+(c_P-X)r a_{-1}-\frac{r^2}{4}a_{-1}^2.
 \label{le:two-modes}
\end{equation}
The square is a complex square, not a squared modulus. Taking absolute
values only after this identity, and putting $x=|a_{-1}|$, proves
\eqref{le:F}. To check attainment in the abstract coefficient ball,
write $(c_P-X)/r=de^{\ii\phi}$ and choose
$a_{-1}=-xe^{\ii\phi}$ and
$a_{-2}=\sqrt{\eta^2-x^2}\,e^{2\ii\phi}$; all three terms then
have the same phase. For $d=0$ any $\phi$ works.
Cauchy--Schwarz bounds the first two terms of its maximum
by $\eta\sqrt{1+d^2}$, proving \eqref{le:cauchy}.
Also $\sqrt{\eta^2-x^2}\le\eta-x^2/(2\eta)$.
Maximizing the resulting concave quadratic over the real line proves
\eqref{le:parabolic}. When $d=0$, set $y=x^2$ and maximize
$\sqrt{\eta^2-y}+y/4$ on $[0,\eta^2]$.
Its maximum is at $y=0$ for $\eta\le2$ and at $y=\eta^2-4$ otherwise,
which proves \eqref{le:centered}.
\end{proof}

The same source and geometry previously give the convenient bound
\begin{equation}
 B_{\rm old}=r^2\{\lambda\sqrt{1+d^2}+\lambda^2/4\}.
 \label{le:old}
\end{equation}
To see this directly, apply $|\ell|\le\lambda$ to the linear contour
moment, use the angular Cauchy--Schwarz bound
$\langle|re^{\ii\theta}+X-c_P|\rangle_\theta
\le r\sqrt{1+d^2}$, and retain the complete center-square payment.
For $q>0$, \eqref{le:cauchy} is strictly below \eqref{le:old}, because
$\eta<\lambda$. This is a strict improvement of that specified complete
allowance; it is not a comparison with an optimal direct estimator.
At $q=1/2,d=0$ the new allowance is $r^2\log(3)/2$, whereas
\eqref{le:old} is $r^2(\log2+(\log2)^2/4)$.
The ratio is less than $17/25$: more than $32\%$ of this full allowance
is saved. For an exact check, the positive $\operatorname{arctanh}$
series gives $\log2>56/81>69/100$ and
\[
 \tfrac12\log3
 =\operatorname{arctanh}(1/2)
 <\frac12+\frac1{24}+\frac1{160}
       +\frac1{7}\sum_{n=3}^{\infty}2^{-(2n+1)}
 =\frac12+\frac1{24}+\frac1{160}+\frac1{672}<\frac{11}{20}.
\]
Meanwhile $(17/25)(69/100+(69/100)^2/4)=0.550137>11/20$.
This finite constant example is not a canonical error sample.

\subsection{What the improved allowance does and does not protect}

For real conjugation-symmetric fields and a fixed counted circle, a
validated lower bound
\begin{equation}
 D_P(u)-r^2F(\eta(u),d(u))>0\quad\hbox{for every }0\le u\le U
 \label{le:margin}
\end{equation}
protects the enclosed real pair. Every error used to acquire $q$, the
count, $c_P$ and $D_P$ must be included. An endpoint instance alone
does not supply \eqref{le:margin}. Nor may the centered subcase be used
by silently moving the fixed circle to $c_P(u)$; the offset $d(u)$ pays
that distinction on the original contour.

If $H/P$ extends as a nonvanishing holomorphic function through the
interior, its logarithm is holomorphic there and both negative modes
vanish. The exact gap error is zero, even though a separate bound on the
joint second-derivative quotient need not vanish for a spatially varying
multiplier. In the full signed-time readout the center and heat-defect
terms supply the cancellation. Integrating that full identity gives
exactly the difference of the endpoint functional
\eqref{prog:centered-error}. Thus identical admissible source paths give
identical exact ranges for full transport and direct endpoint evaluation.
This is a scoped representation equality. Time-correlated data or a
useful heat acquisition may sharpen either calculation; no cost or
universal no-go claim follows.

For the exact control $H=P=(z-X)^2-2(t-T)$, $q=0$ and all new allowances
vanish, yet $D_P=2(t-T)$ reaches zero at its collision. The strict margin
\eqref{le:margin} fails. The remaining controls in Section~\ref{sec:program}
retain the generic heat/logarithm estimates with their own data. Changed
theta coefficients prevent importing \eqref{le:source-value} unchanged;
changed observation kernels fail the physical heat equation. Naming these
source failures does not establish a positive arithmetic margin. The
allowance improvement itself is generic and does not distinguish theta
from a colliding control.

\subsection{A finite arithmetic comparison inside known coverage}
\label{le:pilot}

The bounded pilot fixes
\begin{equation}
 X=140101/10,\quad A=2,\quad h=1,\quad r=9/10,\quad
 N=33,\quad 9/100\le t\le1/10.
 \label{le:pilot-cell}
\end{equation}
All fields and errors are multiplied by the same positive constant
$\exp(\pi X/8)$, independent of $z,t$. It changes neither the physical
heat equation nor the relative error or divisor. The cutoff inequalities
are verified on the entire rectangle. The final calculation uses
192-bit Arb/Acb arithmetic, 64 closed time cells and 512 closed angular
arcs. It proves $\Re(P_N(z)/(z-X)^2)>0$ on the count circle and the
annular boundary circles of radii $4/5$ and $1$, after the common scaling.
Thus the reference has count two and a single annular logarithm throughout
the time interval. The full source budget \eqref{le:source-value} is
included, and $q$ is the maximum of the local source/denominator ratios
on the \emph{same} arcs. Independent global suprema give a weaker bound
and are not used on one side of the comparison alone.

The reference moments are evaluated by the 512-node Laurent trapezoid
rule. If $M_-,M_+$ bound the annular logarithm on $r_-=4/5,r_+=1$,
the paid alias allowance for the weighted coefficient of order $k=1,2$ is
\begin{equation}
 M_+ r_+^k\frac{(r/r_+)^{512}}{1-(r/r_+)^{512}}
 +M_- r_-^k\frac{(r_-/r)^{512}}{1-(r_-/r)^{512}}.
 \label{le:aliases}
\end{equation}
This follows by bounding the positive and negative Laurent coefficient
ladders separately. Both complete allowances include these aliases,
interval time variation, reference-center uncertainty and arithmetic.
The largest coefficient alias is below $5.787\cdot10^{-23}$.
No exterior zero sum is truncated; the full field defines the contour
moments.

\begin{center}
\begin{tabular}{@{}lr@{}}
\toprule
Quantity on the entire pilot interval & Outward summary \\
\midrule
Complete relative-error allowance $q$ & $<0.166454$ \\
Reference discriminant lower bound & $>0.1418896$ \\
E10 transfer allowance \eqref{le:old} & $<0.1545636$ \\
Refined transfer allowance & $<0.1364749$ \\
Worst E10 complete lower bound (approximately) & $-0.0126739660$ \\
Refined complete all-time lower bound & $>0.0054147$ \\
\bottomrule
\end{tabular}
\end{center}

The comparison uses identical acquired $q,c_P,D_P$, source data and
numerical payments. The old sufficient lower bound is negative; this
says nothing adverse about the actual zeros. The new bound is positive
on every time cell and so protects this counted pair. Its transfer
allowance is smaller by more than $11\%$ in the worst-cell summary.
The separate centered $q=1/2$ example's $32\%$ is not the percentage of
this arithmetic pilot. The code, complete small record, method and
nonauthor replay are saved in the investigation's numerics, notes and
reviews folders. They constitute internal certificate checking with a
published source theorem as input, not independent specialist review.

This is a finite improvement of the specified direct estimator inside
an already verified canonical region. It gives no new height/time
coverage and no improved bound on the de Bruijn--Newman constant. Full
signed transport equals the same refined direct endpoint consumer on
identical information.

The scalar heat acquisition also has a definite scoped outcome. In the
same scaled units, its paid uniform bound is above $327656$, while the
raw direct contour allowance is below $196588$. The initial supremum is
already above $327394$; the lateral boundary probability at $U=1/100$
is below $7.358\cdot10^{-15}$. Thus this particular separate-supremum
heat estimate loses to the direct envelope. It does not show that the
actual residual is large, or exclude spatially weighted, signed or
correction-aware acquisitions. The two predeclared refinements and the
failed earlier enclosure are retained in the execution record.

% 19 November 2026. Drafted for Edward Baker with OpenAI Codex assistance.
% Lead actual backend/effort unavailable; helpers explicitly Sol6.1/ultra.
% Analytical integration accepted in Review05; numerical integration assessed separately.
\section{Retaining the arithmetic correction on the fixed cell}
\label{sec:corrected-reference}

The correction in Section~6 changes the approximation center. We now
make that change on the same cell as \eqref{le:pilot-cell}, paying the
difference between the published value formula and a holomorphic reference.
This is a source-acquisition refinement. Direct evaluation of the corrected
reference receives the same information and complete allowance.

\subsection{A local holomorphic reference and its derivative bounds}

Put $k=\pi/8$, $w=z/2+kt$ and
\begin{align}
 p(w)&=1-2\bigl(\sqrt{w/(2\pi)}-N\bigr),\notag\\
 C_0(p)&=\frac{e^{\pi\ii(p^2/2+3/8)}-\ii\sqrt2\cos(\pi p/2)}
                   {2\cos(\pi p)},\qquad
 C_0^\sharp(p)=\overline{C_0(\bar p)},\notag\\
 A_N(w)&=\frac{\sqrt w}{2}
       \bigl\{(w+\ii)C_0(p(w))+(w-\ii)C_0^\sharp(p(w))\bigr\},\notag\\
 J_N(t,z)&=(-1)^{N+1}\frac{\sqrt\pi}{4}e^{-kz-k^2t},\qquad
 C_{\rm hol}(t,z)=J_N(t,z)A_N(w),\qquad Q_N=P_N-C_{\rm hol}.
 \label{cr:reference}
\end{align}
The square roots are principal. All assertions in this section concern
$N=33$, $|\Re z-14010.1|\le2$, $|\Im z|\le1$ and
$9/100\le t\le1/10$; no global continuation of the displayed quotient
is needed.

\begin{proposition}[Explicit correction bounds]
\label{cr:bounds}
On this rectangle the reference is holomorphic in $z$ and respects
conjugation. Its slow profile satisfies
\begin{equation}
 |A_N|<1060000,\qquad |A_N'|<24000,\qquad |A_N''|<1400.
 \label{cr:profile-bounds}
\end{equation}
\end{proposition}

\begin{proof}
The rational bounds $157/50<\pi<22/7$ give
$7004<\Re w<7006.09$, $|\Im w|\le1/2$ and $|w|<7007$.
Writing $\sqrt{w/(2\pi)}=a+\ii b$ and using
$a^2=(|w|+\Re w)/(4\pi)$ and $2ab=\Im w/(2\pi)$ gives
\begin{equation}
 \frac9{50}<\Re p<\frac6{25},\qquad |\Im p|<\frac1{400},
 \qquad |p|<\frac14,\qquad |\cos(\pi p)|>\frac7{10}.
 \label{cr:p-domain}
\end{equation}
The last inequality follows from
$|\cos(u+\ii v)|^2=\cos^2u+\sinh^2v$ and
$\cos(6\pi/25)>7/10$. Thus all denominators and branches remain
clear on a neighborhood of the closed rectangle.

For completeness, set $n=e^{\pi\ii(p^2/2+3/8)}-
\ii\sqrt2\cos(\pi p/2)$ and $d=2\cos(\pi p)$.
On \eqref{cr:p-domain}, the exponential has modulus less than
$1003/1000$ and $\cosh(\pi|\Im p|)<1001/1000$. Using
$|\sin z|\le|z|\cosh|\Im z|$ and $\sqrt2<10/7$ gives
$|n|<5/2$, $|n'|<7/4$, $|n''|<15/2$,
$|d'|<5$, $|d''|<20$ and $|d|>7/5$. Thus
$|C_0|<9/5$, $|C_0'|<8$ and $|C_0''|<90$;
the quotient identities are
\[
 C_0'=\frac{n'-d'C_0}{d},\qquad
 C_0''=\frac{n''-2d'C_0'-d''C_0}{d},\qquad d''=-\pi^2d.
\]
The same bounds hold for $C_0^\sharp$. Write $A_N=gB$, where
$g=\sqrt w/2$ and $B$ is the expression in braces in
\eqref{cr:reference}. Then
\[
 |g|<42,\quad |g'|<\frac1{4\cdot83},\quad
 |g''|<\frac1{8\cdot83^3},\quad
 |p'|<\frac1{200},\quad |p''|<\frac1{2800000}.
\]
Since $|w\pm\ii|<7008$, product and chain rules give
\begin{align*}
 |B|&<2\cdot7008\cdot\frac95,\\
 |B'|&<\frac{18}5+2\cdot7008\cdot\frac8{200},\\
 |B''|&<\frac{32}{200}
   +2\cdot7008\left(\frac{90}{200^2}
                        +\frac8{2800000}\right).
\end{align*}
Substitution into $A_N'=g'B+gB'$ and
$A_N''=g''B+2g'B'+gB''$ proves \eqref{cr:profile-bounds}.
Finally $p(\bar w)=\overline{p(w)}$ and the paired sharp terms give
$A_N(\bar w)=\overline{A_N(w)}$. This also proves the claimed
physical conjugation symmetry.
\end{proof}

Let $T'=x/2+kt$ and let $C_{\rm val}$ denote the correction in
Corollary~6.4 of \cite{Polymath2019}. Substituting the principal-branch
formula for $\mathfrak M_0$ and the source's phase
$U(T')=\exp[-\ii\{(T'/2)\log(T'/(2\pi))-T'/2-\pi/8\}]$
gives exactly
\begin{equation}
 \mathfrak M_0(\ii T')U(T')e^{\pi\ii/8}
 =-\frac{\sqrt\pi}{8}e^{-\pi T'/4}\sqrt{T'}(T'+\ii),
 \qquad C_{\rm val}(t,x+\ii y)=J_N(t,x+\ii y)A_N(T').
 \label{cr:phase}
\end{equation}
Consequently the real-axis corrections agree. Off-axis their difference
is the vertical integral of $A_N'$, which is paid below. Differentiating
the holomorphic reference, with $w_t=k$, $w_z=1/2$, gives the exact
cancellations
\begin{equation}
 (\partial_t+\partial_z^2)C_{\rm hol}
       =\frac14J_NA_N''(w),\qquad
 \mathcal R_Q=\mathcal R_N-\frac14J_NA_N''(w).
 \label{cr:defect}
\end{equation}
The $A_N$ and $A_N'$ terms cancel. In backward time the corrected
error satisfies $(H-Q)_u=(H-Q)_{zz}+\mathcal R_Q$.

\subsection{The complete allowance and its final consumer}

Define the positive leading payment
\begin{equation}
 L(t,x)=e^{k^2t}|\mathfrak M_0(\ii T')|
       =\frac{\sqrt\pi}{8}e^{-kx-k^2t}
                              \sqrt{T'}\sqrt{T'^2+1}.
 \label{cr:leading}
\end{equation}
Write $\mathcal V_P$ for \eqref{le:source-value}. Corollary~6.4,
together with the vertical mismatch, supplies the complete bound
\begin{equation}
 |H-Q_N|\le\mathcal V_Q
 :=\mathcal V_P-L+\Delta,\qquad
 \Delta=|J_N(t,z)|\frac{|y|}{2}\,24000.
 \label{cr:source}
\end{equation}
Both summand errors and the entire retained $\delta_-+\delta_+$
remainder remain in $\mathcal V_P-L$. Equations
\eqref{cr:profile-bounds} and \eqref{cr:leading} give
\begin{equation}
 \frac{\Delta}{L}<\frac{24000}{83\cdot7004}<\frac1{24},
 \qquad \frac{|C_{\rm hol}|}{L}
       <\frac{2\cdot1060000}{83\cdot7004}<4.
 \label{cr:saving}
\end{equation}
Thus more than $23/24$ of the \emph{leading omitted payment} is saved.
This fraction is not a percentage improvement of the complete gap bound.

There is also a concrete count transfer. The previous certificate has
$\mathcal V_P/|P_N|<1/6$ on the circle. Hence
$|C_{\rm hol}/P_N|<2/3$, and Rouch\'e transfers the count two to $Q_N$.
Pointwise, putting $q_P=\mathcal V_P/|P_N|$, one obtains
\begin{equation}
 \frac{\mathcal V_Q}{|Q_N|}
 \le\frac{\mathcal V_P-(23/24)L}{|P_N|-4L}<q_P,
 \quad\text{since }4q_P<23/24.
 \label{cr:relative}
\end{equation}
These inequalities do not supply $Q_N$'s center or discriminant. A final
comparison must acquire them and evaluate
$D_Q-r^2F(\operatorname{arctanh}q_Q,|c_Q-X|/r)$ with every numerical
payment. No $P_N$ denominator or center is reused as though it belonged
to $Q_N$. Direct $Q_N$ evaluation receives the same source information.

The collision controls remain consistent. For an exact colliding
reference the error may vanish while its discriminant reaches zero;
the strict final margin then fails. Modified theta coefficients cannot
import \eqref{cr:source} unchanged, and a changed observation kernel
fails the physical heat equation. These premise failures do not assert
that every collision-free box of a control must fail. The present local
branch and cutoff conditions do not establish an all-height family,
a passage through cutoff changes, or new canonical coverage.

\subsection{The complete corrected comparison}

The second calculation keeps exactly the previous $192$-bit arithmetic,
$64$ closed time cells, $512$ closed arcs and three radii. It evaluates
$Q_{33}$ directly and reacquires its count, denominator and annular
logarithm from $\Re(Q_{33}(z)/(z-X)^2)>0$ on every boundary arc.
The two logarithmic moments use the same fully paid Laurent alias formula
\eqref{le:aliases}; interval time variation is included. Both term errors,
the retained remainder and $12000|J_N||y|$ enter the local-arc ratio.
The old $P_{33}$ clearance and moments are not substituted for $Q_{33}$'s.

\begin{center}
\begin{tabular}{lr}
\toprule
Quantity over the entire fixed interval & Outward summary\\
\midrule
Corrected relative-error allowance $q_Q$ & $<0.015234762$\\
Corrected reference discriminant $D_Q$ & $>0.1617006119$\\
Refined corrected transfer allowance & $<0.012369746$\\
E10 complete lower bound at $Q$ & $>0.1491891639$\\
Refined complete lower bound at $Q$ & $>0.1493308659$\\
Previous refined complete lower bound at $P$ & $>0.0054147305$\\
\bottomrule
\end{tabular}
\end{center}

The refined corrected lower bound exceeds the previous evaluated lower
bound in every one of the $64$ closed time cells. In particular
\begin{equation}
 D_H(t)>0.1493308659\qquad(9/100\le t\le1/10)
 \label{cr:pilot-margin}
\end{equation}
on the same counted circle. The real pair has gap greater than $0.7728$,
since $(0.7728)^2/4=0.14930496<0.1493308659$.
This is a stronger sufficient lower bound for the same canonical pair;
the zeros have not changed. The gain uses retained correction information,
and is shared by direct evaluation with the same information. It does not
show an advantage of a time-integrated representation.

The scaled complete source allowance is below $21410.476$ and the
corrected denominator exceeds $941112.53$. Their independent ratio is
not used in place of the sharper local-arc ratios. The complete corrected
defect is conservatively bounded by $|\mathcal R_N|+350|J_N|$;
its scaled contour bound is below $357.797450$. This is an upper
triangle estimate, not an evaluation of the signed corrected defect.
No second scalar heat-propagation comparison is claimed.
The code, full record, separate assessment and sole scratch replay are
saved in the investigation. They retain the imported source theorem and
internal LLM-review limits. The fixed cell remains inside known coverage;
no new height/time region or improved de Bruijn--Newman bound follows.

% 19 November 2026. Drafted for Edward Baker with OpenAI Codex assistance.
% Root settings unavailable; B03 Astra/ultra and E14 Sol6.1/ultra dispatches.
% Integration of the E14 proof; assessment status is recorded in Review07.
\section{Full-phase correction and the pointwise strip consumer}
\label{sec:uniform-correction}

The fixed-cell gain raises two different questions: whether the correction
can be controlled through the whole Riemann--Siegel phase range, and whether
the resulting approximation proves a positive final margin. We resolve
the first and compare complete pointwise allowances for the second.
The comparison retains an explicit unacquired arithmetic lower bound;
it proves no new zero-free region.

\subsection{One family and removable apparent poles}

Keep \eqref{cr:reference}, now with the natural index and the family
\begin{equation}
 \frac9{100}\le t\le\frac1{10},\qquad x\ge14000,\qquad |y|\le1,
 \quad W=\frac{x}{2}+\frac{\pi t}{8},\quad
 w=W+\frac{\ii y}{2},\quad
 N=\left\lfloor\sqrt{\frac W{2\pi}}\right\rfloor.
 \label{uc:family}
\end{equation}
For clarity distinguish the published value correction and its lift:
\begin{align}
 C_{\rm val,N}&=J_N A_N(W),& Q_{\rm val,N}&=P_N-C_{\rm val,N},\\
 C_{\rm hol,N}&=J_N A_N(w),& Q_{\rm hol,N}&=P_N-C_{\rm hol,N}.
 \label{uc:two-references}
\end{align}
Only the second reference with a \emph{fixed} index is holomorphic.
The natural-index patchwork is not asserted to be holomorphic across
cutoff seams. The value reference is sufficient for a pointwise
nonvanishing test, which takes no spatial derivative.

\begin{lemma}[Full-phase analytic bounds]
The functions $C_0$ and $C_0^\sharp$ in \eqref{cr:reference} extend to
entire even functions. On
$|\Re p|\le101/100$, $|\Im p|\le1/100$ they satisfy
\begin{equation}
 |C_0|<4,\qquad |C_0'|<20,\qquad |C_0''|<200,
 \label{uc:C-bounds}
\end{equation}
and the same bounds hold for the sharp function. Every natural phase
and its vertical continuation in \eqref{uc:family} lies in this patch.
\end{lemma}

\begin{proof}
At $p=j+1/2$, the exponential in the numerator is
$\ii(-1)^{j(j+1)/2}$, while
$\sqrt2\cos(\pi j/2+\pi/4)=(-1)^{j(j+1)/2}$.
The numerator therefore vanishes at every simple zero of the
denominator. Removal of these singularities proves entire extension;
evenness is inherited from the original quotient.

On the boundary of $|\Re p|\le5/4$, $|\Im p|\le1/4$, the denominator
has modulus at least $\sqrt2$: this follows on vertical sides from
$|\cos(a+\ii b)|^2=\cos^2a+\sinh^2b$, and on horizontal sides from
$2\sinh(\pi/4)>\sqrt2$. The numerator has modulus below
$e^{5\pi/16}+\sqrt2\cosh(\pi/8)<5$. Hence the quotient is below
$5/\sqrt2<4$, including the interior by the maximum principle.
Cauchy disks of radius $1/5$ centered on the smaller rectangle remain
inside the larger one, proving the derivative bounds.

On the vertical segment through $W\ge7000$,
$|p_N'|\le1/\sqrt{2\pi W}$. Its phase displacement from the real
phase is at most $1/(2\sqrt{2\pi W})<1/100$.
The real natural phase lies in $(-1,1]$, with $-1$ included by
one-sided limits. This proves the asserted full phase coverage, without
excluding neighborhoods of $\pm1/2$.
\end{proof}

\subsection{Derivative, mismatch and complete source payments}

\begin{proposition}[Uniform correction payments]
Throughout \eqref{uc:family}, with primes denoting derivatives in $w$,
\begin{equation}
 |A_N|<5W^{3/2},\qquad |A_N'|<9W,\qquad
 |A_N''|<40\sqrt W.
 \label{uc:A-bounds}
\end{equation}
For fixed $N$, the complete heat defect is
\begin{equation}
 R_{Q_{\rm hol,N}}=R_{P_N}-\frac{J_N A_N''}{4},\qquad
 |R_{Q_{\rm hol,N}}|\le |R_{P_N}|+10\sqrt W\,|J_N|.
 \label{uc:defect}
\end{equation}
The value-to-lift mismatch is bounded by
\begin{equation}
 |C_{\rm hol,N}-C_{\rm val,N}|\le M_N
 :=\frac92 W|y||J_N|,\qquad
 \frac{M_N}{L}\le\frac{9|y|}{\sqrt W}\le\frac9{\sqrt{7000}}<\frac19,
 \quad L=\frac{|J_N|}{2}\sqrt W\sqrt{W^2+1}.
 \label{uc:mismatch}
\end{equation}
\end{proposition}

\begin{proof}
Put $B=((w+\ii)C_0+(w-\ii)C_0^\sharp)/2$, so $A_N=\sqrt w B$.
Writing $d=(2\pi W)^{-1/2}$, the full product and chain rules give
\begin{align*}
 |B|&\le4(W+3/2),\\
 |B'|&\le4+20(W+3/2)d,\\
 |B''|&\le40d+(W+3/2)\{200d^2+10d/W\}.
\end{align*}
Use $|\sqrt w|\le\sqrt{W+1/2}$,
$|(\sqrt w)'|\le1/(2\sqrt W)$ and
$|(\sqrt w)''|\le1/(4W^{3/2})$.
Substitution into $A_N'=(\sqrt w)'B+\sqrt w B'$ and
$A_N''=(\sqrt w)''B+2(\sqrt w)'B'+\sqrt w B''$
gives \eqref{uc:A-bounds}; for explicit overpayments one can use
$\sqrt W>80$, $d<(5/12)/\sqrt W$,
$W+3/2<(1001/1000)W$ and
$\sqrt{W+1/2}<(1001/1000)\sqrt W$.
All terms, including $p_N''$, are retained in the displayed bounds.

The $A_N$ and $A_N'$ terms cancel in $\partial_t+\partial_{zz}$
exactly as in \eqref{cr:defect}, giving \eqref{uc:defect}.
Integration along the vertical segment gives the mismatch.
Finally $W^2+1\ge W^2$ and $7000>81^2$ prove the last comparison.
This is a triangle defect bound, with no assumed signed cancellation.
\end{proof}

Let $\mathcal V_N=S_N+L(\delta_-+\delta_+)$ denote the complete
retained source budget from \eqref{cr:source}, using the natural $N$ in
both summand-error sums and the full retained remainder. The source
estimate of \cite[Corollary 6.4]{Polymath2019} and
\eqref{uc:mismatch} give
\begin{align}
 |H-P_N|&\le V_{\rm om}:=\mathcal V_N+L,\notag\\
 |H-Q_{\rm val,N}|&\le V_{\rm val}:=\mathcal V_N,\notag\\
 |H-Q_{\rm hol,N}|&\le V_{\rm hol}:=\mathcal V_N+M_N.
 \label{uc:budgets}
\end{align}
In particular $V_{\rm om}-V_{\rm hol}>8L/9$. This saves a fraction
of the omitted leading payment, not of the entire error or final margin.

\subsection{A quantitative comparison of complete pointwise allowances}

\begin{theorem}[Retained-value improvement and the direct comparator]
On \eqref{uc:family} define
\[
 m_{\rm om}=|P_N|-V_{\rm om},\qquad
 m_{\rm val}=|Q_{\rm val,N}|-V_{\rm val},\qquad
 m_{\rm hol}=|Q_{\rm hol,N}|-V_{\rm hol}.
\]
These are lower bounds for $|H|$, and
\begin{equation}
 \boxed{m_{\rm val}>m_{\rm om}+\frac3{40}L,
 \qquad m_{\rm hol}\le m_{\rm val}.}
 \label{uc:complete-comparison}
\end{equation}
Both comparisons remain true after division by the same positive
carrier modulus. Neither asserts that $m_{\rm val}$ is positive.
\end{theorem}

\begin{proof}
For real $p\in[-1,1]$, \cite[(58)]{Polymath2019} gives $|C_0(p)|\le1/2$.
There is also the elementary bound
$0<\Re C_0(p)\le\cos(\pi/8)/2$. To prove it, put
$q=|p|$, $a=q-1/2$ and $b=(q+1/2)/2$. Cancellation in the real part gives
\[
 \Re C_0(p)=\frac{\sin(\pi ab)}{2\sin(\pi a)}.
\]
The expression is interpreted continuously at $a=0$.
For $u=\pi|a|\in[0,\pi/2]$ and $b\in[1/4,3/4]$,
$\sin(bu)/\sin u$ increases in both arguments. Its $u$ derivative has
the sign of $b\cot(bu)-\cot u$, which is positive because $v\cot v$
decreases on $(0,\pi/2]$. The maximum is therefore
$\sin(3\pi/8)=\cos(\pi/8)$.

Since $A_N(W)=\sqrt W\{W\Re C_0-\Im C_0\}$,
\[
 \frac{|C_{\rm val,N}|}{L}
 \le\frac{W\cos(\pi/8)+1}{\sqrt{W^2+1}}
 \le\cos(\pi/8)+\frac1W<\frac{37}{40}.
\]
For the last rational bound use $\cos(\pi/8)<231/250$ and
$1/W\le1/7000<1/1000$. The reverse triangle inequality now gives
$m_{\rm val}-m_{\rm om}\ge L-|C_{\rm val,N}|>3L/40$.
Finally $Q_{\rm hol,N}-Q_{\rm val,N}=C_{\rm val,N}-C_{\rm hol,N}$,
so $|Q_{\rm hol,N}|\le|Q_{\rm val,N}|+M_N$, proving the second claim.
\end{proof}

The first comparison retains correction information that the omitted
budget discards. The second gives that same information to the direct
consumer. It is restricted to exact pointwise reference values and a
separately paid vertical triangle allowance. It does not compare arbitrary
interval implementations, computational costs, correlated residuals,
counted-pair functionals or heat propagation. In particular it does not
withdraw the fixed-cell discriminant gain of Section~9.

\subsection{Cutoff seams and the remaining arithmetic input}

At $W=2\pi j^2$, equivalently $x=4\pi j^2-\pi t/4$, the old index is
$j-1$ and the new one is $j$. Put
$d_j=p_j(t,s_+)+p_j(t,s_-)$, as in \eqref{as:cutoff-jump}.
Since $J_{j-1}=-J_j$ and the real seam phases are $-1,+1$,
evenness gives the exact signed jumps
\begin{align}
 \Delta Q_{\rm val}&=d_j-2J_jA_j(W),\notag\\
 \Delta Q_{\rm hol}&=d_j-J_j\{A_j(w)+A_{j-1}(w)\}.
 \label{uc:seams}
\end{align}
The two profiles in the second line are not identified off the real axis.
Both one-sided source estimates apply at the seam by continuity; hence
$|\Delta Q|\le V_{\rm old}+V_{\rm new}$ for each type of reference,
with its own complete budgets. A change of reference transfers the error
by $E_{\rm new}=E_{\rm old}-\Delta Q$, including the jump payment.
Pointwise coverage uses the natural half-open cells and their endpoints.
Differentiating the floor, or treating the patchwork as a single heat
solution without seam terms, is not justified by these formulas.

The chosen final consumer is precisely \eqref{as:strip-final} with
\begin{equation}
 t_0=1/10,\quad y_0=1/5,\quad X=14000,\qquad
 x\ge14000+\sqrt{24/25},\quad 1/5\le y\le\sqrt{4/5}.
 \label{uc:consumer}
\end{equation}
The nonvanishing carrier $\mathcal B_{t_0}$ in
\eqref{as:normalizer} permits the normalized statement. What remains
unacquired is an independently proved lower function
\[
 \ell(x,y)\le
 \left|\frac{Q_{\rm val,N}(t_0,x+\ii y)}{\mathcal B_{t_0}(x+\ii y)}\right|,
 \qquad
 \ell(x,y)>\frac{\mathcal V_N(t_0,x+\ii y)}
 {|\mathcal B_{t_0}(x+\ii y)|}
\]
at every point of \eqref{uc:consumer}, with a complete finite-height/tail
split and all cutoff interfaces. A positive unscaled infimum over
unbounded height is not requested. Defining $\ell$ from the desired
positivity would not acquire it.

The family includes the finite interval
$[14000,14000+\sqrt{24/25})$, but acquires no positivity there.
It supplies neither the initial verification \eqref{as:strip-initial}
nor the barrier \eqref{as:strip-barrier} down to time zero. Only if
all three criterion inputs were established would the parameter
arithmetic yield $\Lambda\le1/10+1/50=3/25$; this is not a new bound.

The collision screens retain their precise failed premises. The
Euler-factor, weighted-square, finite-prime and reflected/positive-primitive
controls do not inherit the literal theta source budget. That failure
does not turn a small approximation error into a zero-sensitive margin:
at a zero, any valid complete allowance gives $|Q|-V\le0$.
The quadratic $H=Q=(z-X)^2-2(t-T)$ has zero approximation error and
still collides. Finite-source controls lack the required all-real anchor;
RH-conditional escaping controls keep their hypothesis and location
qualifiers; changed observation kernels pay their actual forcing.
Thus the complete comparison is useful, while final positivity remains
an arithmetic acquisition problem.

% 19 November 2026. Drafted for Edward Baker with OpenAI Codex assistance.
% Root settings unavailable; B04 Astra/ultra and E16 Sol6.1/ultra dispatches.
% Assessment and source hashes are recorded in Review08.
\section{A quantified limitation of a three-prime lower allowance}
\label{sec:arithmetic-obstruction}

We now test an independent arithmetic reduction for the missing strip
margin. The finite Euler multiplier of \cite[Section~8.5]{Polymath2019}
cancels selected coefficients before taking absolute values. A paired
ellipse retains additional dependence between the two principal sums.
Even this stronger reduction has a negative complete allowance on an
included boundary segment. This is a limitation of the specified
estimator, not a zero of the theta field or an exclusion of all arithmetic
methods.

\subsection{Exact finite convolution and the paid allowance}

Fix $t=1/10$ and use the normalized quantities of Section~5:
\begin{align}
 s_0&=(1+y-\ii x)/2,& s&=s_0+\alpha(s_0)/20=\sigma+\ii\tau,
 &b_n&=\exp(\log^2 n/40),\notag\\
 N&=\left\lfloor\sqrt{x/(4\pi)+1/160}\right\rfloor,&
 S&=\sum_{n\le N}b_n n^{-s},&
 T&=\sum_{n\le N}b_n n^{y-s}.
 \label{af:definitions}
\end{align}
Thus $f=P_N/\mathcal B=S+\gamma\sum_{n\le N}b_n
n^{y-\bar s-\kappa}$. Put $c=C_{\rm val,N}/\mathcal B$,
$q=Q_{\rm val,N}/\mathcal B=f-c$ and
$v=\mathcal V_N/|\mathcal B|$, with both full term-error sums and
the entire retained delta remainder in $\mathcal V_N$. In particular
$|H/\mathcal B-q|\le v$.

The chosen multiplier and its exact coefficients are
\begin{align}
 E&=\prod_{p\in\{2,3,5\}}(1-b_p p^{-s})
   =\sum_{d\mid30}\lambda_d d^{-s},&
 \lambda_d&=\prod_{p\mid d}(-b_p),\notag\\
 \beta_m&=\sum_{\substack{d\mid(m,30)\\m/d\le N}}
                  \lambda_d b_{m/d},&
 a_m&=\sum_{\substack{d\mid(m,30)\\m/d\le N}}
                  \lambda_d(m/d)^y b_{m/d}.
 \label{af:convolution}
\end{align}
For $1\le m\le30N$ these are real, with $\beta_1=a_1=1$,
$ES=\sum_m\beta_m m^{-s}$ and $ET=\sum_m a_m m^{-s}$.
The cutoff indicator includes all edge terms beyond $N$; the heat weights
are not assumed multiplicative. Equal-frequency cancellation uses the
actual identity $(dn)^{-\ii\tau}=d^{-\ii\tau}n^{-\ii\tau}$.

Let $g=|\gamma|<1$ and suppose $E\ne0$. The kappa replacement costs
\begin{equation}
 D_\kappa=g\sum_{n\le N}b_n n^{y-\sigma}
                    \bigl(e^{|\kappa|\log n}-1\bigr),
 \qquad |f-(S+\gamma\overline T)|\le D_\kappa.
 \label{af:kappa}
\end{equation}
Define
\begin{align}
 h_m&=\max\left\{|\beta_m-ga_m|,
          \frac{1-g}{1+g}|\beta_m+ga_m|\right\},\notag\\
 K&=1-g-\sum_{m=2}^{30N}h_m m^{-\sigma},\notag\\
 \ell_E&=\frac{K}{|E|}-D_\kappa-|c|,
 &m_E&=\ell_E-v.
 \label{af:allowance}
\end{align}

\begin{lemma}[Paired-ellipse modulus estimate]
The quantities in \eqref{af:allowance} satisfy
$\ell_E\le|q|$ and
\begin{equation}
 m_E\le m_{\rm dir}:=|f-c|-v\le|H/\mathcal B|.
 \label{af:direct}
\end{equation}
These signed inequalities also hold when $K<0$.
\end{lemma}

\begin{proof}
Put $U=ES$, $V=ET$ and $\rho=\gamma E/\bar E=g e^{\ii\theta}$.
Then $E(S+\gamma\bar T)=U+\rho\bar V$.
After rotation by $e^{-\ii\theta/2}$, its $m$th pair is
\[
 m^{-\sigma}\bigl\{(\beta_m+ga_m)\cos\phi_m
                 +\ii(\beta_m-ga_m)\sin\phi_m\bigr\},
 \qquad \phi_m=-\tau\log m-\theta/2.
\]
Use the norm
$\|w\|_g^2=(\Re w/(1+g))^2+(\Im w/(1-g))^2$.
The constant pair has norm one; the other pairs have norm at most
$h_m m^{-\sigma}/(1-g)$. Since $|w|\ge(1-g)\|w\|_g$,
the reverse triangle inequality gives $|U+\rho\bar V|\ge K$.
Pay \eqref{af:kappa}, divide by the exact $|E|$, and pay $|c|$.
This gives $\ell_E\le|f-c|$; the complete source enclosure gives
\eqref{af:direct}. Negative lower allowances cause no sign reversal
because the divisor remains the exact positive $|E|$.
\end{proof}

The ellipse preserves within-index dependence but its triangle sum
discards the relations among different phases $\tau\log m$. Paying
$|c|$ separately is another explicit restriction of this estimator.
Direct evaluation receives the same coefficients, phases and complete
payments, and can use the same ellipse estimate. No computational-cost
or universal interval-implementation comparison follows.

\subsection{Eight surviving primes force a negative allowance}

\begin{theorem}[Scoped three-prime obstruction]
On the boundary segment
\begin{equation}
 \mathcal F_0:\quad t=1/10,\quad y=1/5,\quad
 x_0:=14000+\sqrt{24/25}\le x\le30000,
 \label{af:segment}
\end{equation}
the exact allowance \eqref{af:allowance} satisfies
\begin{equation}
 \boxed{\ell_E<-1/400,\qquad m_E<-1/400.}
 \label{af:negative}
\end{equation}
Both one-sided cutoff references satisfy the same conclusion at each
seam in this segment. The theorem concerns the fixed multiplier
$\{2,3,5\}$, anchor one, individual paired-ellipse radii and separate
kappa, correction and source payments.
\end{theorem}

\begin{proof}
First $N\ge33$ since $x>14000$ and $\pi<22/7$.
Writing $T_0=x/2$, the exact alpha formula gives
\begin{align*}
 \Re\alpha(3/5-\ii T_0)-\tfrac12\log(T_0/(2\pi))
 &=\frac{3/10}{T_0^2+9/25}-\frac{2/5}{T_0^2+4/25}
      +\tfrac14\log\bigl(1+9/(25T_0^2)\bigr)\\
 &\le\frac{39/100}{T_0^2}-\frac{2/5}{T_0^2+4/25}<0.
\end{align*}
Also $x/(4\pi)<2500<e^8$, using $\pi>3$ and $e>8/3$.
Hence $\sigma<3/5+8/40=4/5$. A lower bound
\[
 \Re\alpha(3/5-\ii T_0)>
 \tfrac12\log(T_0/(2\pi))-\frac{2}{5T_0^2}>0
\]
gives $\sigma>3/5$. Here $T_0>7000$ and $T_0/(2\pi)>1000$.
The general source bound \cite[(20)]{Polymath2019} yields
\[
 0<g\le e^{1/250}\bigl(x/(4\pi)\bigr)^{-1/10}<3/5:
 \qquad 1000^{1/10}>9/5,\quad e^{1/250}<27/25.
\]
The last estimates follow from $(9/5)^{10}<1000$ and
$e^u\le(1-u)^{-1}$ for $0<u<1$. For $p=2,3,5$,
$0<\log p<2$ and $\sigma>3/5$ imply
$b_p p^{-\sigma}<1$. Thus $0<|E|<8$.

Let $\mathcal P=\{7,11,13,17,19,23,29,31\}$ and
$P_j(u)=\sum_{h=0}^j u^h/h!$. The following rational table provides
$\log p>r$, $r/5+r^2/40>v_p$ and $P_4(v_p)/p>k_p/1000$:
\begin{center}
\begin{tabular}{rrrrr}
\toprule
$p$&$r$&$v_p$&$k_p$&$(a,j)$\\
\midrule
7&$19/10$&$47/100$&228&$(2,1)$\\
11&$23/10$&$59/100$&163&$(3,2)$\\
13&$5/2$&$13/20$&147&$(3,2)$\\
17&$14/5$&$3/4$&124&$(3,1)$\\
19&$29/10$&$79/100$&115&$(3,1)$\\
23&$31/10$&$43/50$&102&$(4,3)$\\
29&$33/10$&$93/100$&87&$(4,2)$\\
31&$17/5$&$193/200$&84&$(4,2)$\\
\bottomrule
\end{tabular}
\end{center}
For verification, $e<P_6(1)+1/4410<68/25$, and each row satisfies
$(68/25)^a/P_j(a-r)<p$, proving $e^r<p$.
The remaining columns are rational polynomial inequalities.
Consequently
\begin{equation}
 \sum_{p\in\mathcal P} b_p p^{-4/5}
 =\sum_{p\in\mathcal P}p^{-1}
       e^{\log p/5+\log^2p/40}
 >\frac{\sum_p k_p}{1000}=\frac{21}{20}.
 \label{af:prime-mass}
\end{equation}
All eight primes lie below $N$ and are coprime to 30, so
$\beta_p=b_p$ and $a_p=p^{1/5}b_p\ge b_p$. The second branch of the
maximum therefore gives $h_p\ge(1-g)b_p$. Using $\sigma<4/5$,
\[
 K<1-g-\frac{21}{20}(1-g)
   =-\frac{1-g}{20}<-\frac1{50}.
\]
The already-defined signed quotient then satisfies
$K/|E|<-1/(50|E|)<-1/400$. The three remaining payments are
nonnegative, proving \eqref{af:negative}. The upper bound eight is
used only to bound a negative allowance from above, not to replace
its denominator in a lower-bound formula.
\end{proof}

\subsection{Coverage, controls and the next missing information}

At $x_j=4\pi j^2-\pi/40$, apply \eqref{af:convolution} separately
with indices $j-1$ and $j$. Both are at least 33 at every seam in
\eqref{af:segment}; the same eight-prime witness applies.
The old endpoint is a one-sided limit and the new index owns the
natural endpoint. Transferring a fixed reference still requires the
signed jump and full error payment in \eqref{uc:seams}.

This failure includes the required strip's lower endpoint. Moving a
finite-height/tail join cannot make this estimator certify the whole
strip. An effective tail onset and paid overlap remain independent
unacquired inputs to any successful replacement; eventual classical
two-carrier coverage is not new. The initial verification and full
intermediate barrier remain separate obligations as in
\eqref{uc:consumer}. No tail computation or phase scan supports this
analytical obstruction.

The collision controls distinguish literal theta arithmetic from
selected multiplicative identities. Finite-prime and finite-source
perturbations can preserve some identities while changing the full
coefficients, reflection and source remainder. Shape, primitive and
reflected controls do not supply that complete source premise; conditional
escaping controls retain their hypothesis and location qualifications.
Changed observation kernels pay their actual forcing. At any included
zero, a valid complete enclosure forces every valid $\ell-v\le0$.
For example $H=Q=(z-x_c)^2-2(t-1/4)$, $\mathcal B=1$, $v=0$ has
the zero $z=x_c+\ii\sqrt{3/10}$ at $t=1/10$, within the selected
strip when $x_c\ge x_0$. It has perfect approximation but does not
have theta's arithmetic.

The theorem excludes neither larger or adaptive multipliers nor signed
groups, height-phase constraints, joint correction/remainder estimates
or contour methods. The specific next acquisition is control of a
phase-preserving group by an independent quantitative input. Simply
shrinking correction constants cannot repair a cost already exceeding
the anchor before those payments are included. No positive strip floor,
new zero-free region or improved Newman bound is obtained here.

% 19 November 2026. Drafted for Edward Baker with OpenAI Codex assistance.
% E17 helper explicitly dispatched gpt-6.1-sol/ultra; backend unattested.
\section{A common dyadic phase still leaves an envelope obstruction}
\label{sec:dyadic-group-obstruction}

We strengthen Section~\ref{sec:arithmetic-obstruction} by grouping all
frequencies with the same odd core and retaining their common dyadic
phase across groups. The remaining full-circle envelope still has a
negative complete allowance. All source quantities $q=f-c$, $v$,
$D_\kappa$, $E$, $\beta_m,a_m$, $h_m$ and the segment $\mathcal F_0$
are exactly those of \eqref{af:definitions}--\eqref{af:segment}.
In particular, both full term-error sums and both retained delta payments
remain in $v$, and $|H/\mathcal B-q|\le v$. No new source estimate is
imported.

\subsection{Exact dyadic groups and their allowance}

Write $\rho=\gamma E/\bar E=g e^{\ii\theta}$ as before. Let
$\mathcal R_N$ be the positive odd integers at most $30N$, and put
\[
 I_1=\{j\ge1:2^j\le30N\},\qquad
 I_r=\{j\ge0:r2^j\le30N\}\quad(r\ge3\text{ odd}).
\]
For real auxiliary phases define the finite polynomials and groups
\begin{align}
 B_r(u)&=\sum_{j\in I_r}\beta_{r2^j}2^{-j\sigma}e^{\ii ju},&
 A_r(u)&=\sum_{j\in I_r}a_{r2^j}2^{-j\sigma}e^{\ii ju},\notag\\
 G_r(\eta,u)&=r^{-\sigma}
    \bigl(e^{\ii\eta}B_r(u)+g e^{-\ii\eta}\overline{A_r(u)}\bigr).
 \label{dg:groups}
\end{align}
With $F=S+\gamma\bar T$, the exact finite convolution gives
\begin{equation}
 e^{-\ii\theta/2}EF=Z_0+
 \sum_{r\in\mathcal R_N}G_r(\eta_r,u),\qquad
 Z_0=e^{-\ii\theta/2}+g e^{\ii\theta/2},\quad
 u=-\tau\log2,\quad\eta_r=-\tau\log r-\theta/2.
 \label{dg:identity}
\end{equation}
Every $m\ge2$ has its unique dyadic representation; the excluded
$r=1,j=0$ term is the displayed anchor. All cutoff terms through $30N$
are retained, without assuming multiplicativity of $b_n$.

Use $\|w\|_g^2=(\Re w/(1+g))^2+(\Im w/(1-g))^2$.
Its anchor norm is one and $|w|\ge(1-g)\|w\|_g$. Set
\begin{align}
 R_r(u)&=(1-g)\sup_{\eta\in[0,2\pi]}\|G_r(\eta,u)\|_g,\notag\\
 C_{\rm dyad}&=\sup_{u\in[0,2\pi]}\sum_rR_r(u),&
 C_{\rm sep}&=\sum_r\sup_{u\in[0,2\pi]}R_r(u),\notag\\
 K_{\rm dyad}&=1-g-C_{\rm dyad},&
 \ell_{\rm dyad}&=K_{\rm dyad}/|E|-D_\kappa-|c|,\qquad
 m_{\rm dyad}=\ell_{\rm dyad}-v.
 \label{dg:allowance}
\end{align}
The coefficients, $N,\sigma,g$ are frozen at the actual observation
while these auxiliary suprema are taken. Only the core rotations are
independent; the same $u$ is retained across all groups.

For a real matrix, $\|\cdot\|_2$ denotes its Euclidean operator norm.
Direct expansion of $G_r$ in $\cos\eta,\sin\eta$ gives
\begin{equation}
 M_r(u)=\begin{pmatrix}
 \dfrac{\Re B_r+g\Re A_r}{1+g}&
       -\dfrac{\Im B_r+g\Im A_r}{1+g}\\[4pt]
 \dfrac{\Im B_r-g\Im A_r}{1-g}&
        \dfrac{\Re B_r-g\Re A_r}{1-g}
 \end{pmatrix},\qquad
 R_r(u)=(1-g)r^{-\sigma}\|M_r(u)\|_2.
 \label{dg:matrix}
\end{equation}
Indeed the scaled real coordinates are
$r^{-\sigma}M_r(u)(\cos\eta,\sin\eta)^{\mathsf T}$, whose supremum
is exactly that operator norm. Reverse triangle in \eqref{dg:identity}
therefore gives $|EF|\ge1-g-\sum_rR_r(u)\ge K_{\rm dyad}$.
Paying \eqref{af:kappa} and $|c|$ proves $\ell_{\rm dyad}\le|q|$,
including when $K_{\rm dyad}<0$. The divisor remains the exact $|E|$.

Each $j$th matrix coefficient in \eqref{dg:matrix} is a diagonal
matrix followed by rotation through $ju$, with norm
$2^{-j\sigma}h_{r2^j}/(1-g)$. Its triangle inequality and the unique
dyadic partition imply
\begin{equation}
 C_{\rm dyad}\le C_{\rm sep}\le
 C_{\rm ind}:=\sum_{m=2}^{30N}h_m m^{-\sigma},\qquad
 m_E\le m_{\rm dyad}\le m_{\rm dir}=|f-c|-v.
 \label{dg:comparison}
\end{equation}
Direct evaluation receives every polynomial, actual phase and complete
payment and may use this same grouping. No cost advantage is asserted.

\subsection{A Fourier mean forces the same negative constant}

\begin{theorem}[Scoped common-phase envelope obstruction]
On $\mathcal F_0$ in \eqref{af:segment}, including both one-sided cutoff
references, $\ell_{\rm dyad}<-1/400$ and $m_{\rm dyad}<-1/400$.
The scope is the fixed three-prime multiplier, anchor one, odd-core
dyadic groups, independent core rotations, one common full-circle
dyadic-phase envelope, and separate kappa, correction and source costs.
\end{theorem}

\begin{proof}
Write $\langle\cdot\rangle=(2\pi)^{-1}\int_0^{2\pi}(\cdot)\,du$.
For odd $r\ge3$, the real constant coefficients of the finite
polynomials give
\[
 \langle M_r\rangle=\operatorname{diag}
 \left(\frac{\beta_r+ga_r}{1+g},
       \frac{\beta_r-ga_r}{1-g}\right).
\]
The integral triangle inequality, applied to $M_r$ on each unit vector,
proves $\|\langle M_r\rangle\|_2\le\langle\|M_r\|_2\rangle$.
Thus \eqref{dg:matrix} implies
\[
 \langle R_r\rangle\ge h_r r^{-\sigma},\qquad
 C_{\rm dyad}\ge\left\langle\sum_rR_r\right\rangle
 \ge\sum_{\substack{r\ge3\ {\rm odd}\\r\le30N}}h_r r^{-\sigma}.
\]
The core $r=1$ has zero mean matrix and nonnegative cost. No limiting
argument or nonvanishing assumption enters these finite norm identities.

Section~\ref{sec:arithmetic-obstruction} proves uniformly on
$\mathcal F_0$ that $N\ge33$, $3/5<\sigma<4/5$, $0<g<3/5$,
$0<|E|<8$, and the eight-prime mass \eqref{af:prime-mass} exceeds
$21/20$. Each prime there is an odd core with
$\beta_p=b_p$, $a_p=p^{1/5}b_p$, hence $h_p\ge(1-g)b_p$.
Consequently
\[
 C_{\rm dyad}>\frac{21}{20}(1-g),\qquad
 K_{\rm dyad}<-\frac{1-g}{20}<-\frac1{50},\qquad
 \frac{K_{\rm dyad}}{|E|}<-\frac1{50|E|}<-\frac1{400}.
\]
The three remaining payments are nonnegative. The bound eight is used
only to upper-bound this already-defined negative allowance.
\end{proof}

For this estimator a positive complete allowance would require exactly
\[
 C_{\rm dyad}<1-g-|E|(D_\kappa+|c|+v)\le1-g,
\]
which the mean bound prevents. Clipping $\ell_{\rm dyad}$ to zero
cannot beat $v\ge0$ strictly. These statements upper-bound a lower
allowance, not $|q|$ or $|H/\mathcal B|$.
The smaller actual-phase cost $\sum_rR_r(u(x))$ also gives a valid
allowance, whose sign is undetermined here. A circle mean is not a
pointwise lower bound at $u(x)$, and no adverse maximizing phases are
claimed to occur at one physical height.

\subsection{Interfaces and the information still missing}

The bounds $25/8<\pi<22/7$ and $x_0<14001$ give
$x_{33}<14000<x_0<x_{34}$ and $x_{48}<30000<x_{49}$, where
$x_j=4\pi j^2-\pi/40$. Thus the natural indices are $33,\ldots,48$,
with internal seams $j=34,\ldots,48$. Apply each finite formula on
its half-open natural cell and old one-sided limit; the new index owns
the seam and $30000$ is included with $N=48$. Both references retain
all eight primes. A fixed-reference transfer still pays the exact
signed jump and full error in \eqref{uc:seams}; no floor derivative
or unpaid gluing is used.

Failure at the included $x_0$ cannot be repaired by moving a tail join.
The rest of the two-dimensional finite strip, effective tail onset and
overlap, initial verification and full intermediate barrier remain
unacquired. The collision screen in Section~\ref{sec:arithmetic-obstruction}
also applies: selected arithmetic or dyadic identities do not supply the
complete literal theta source or a strict margin. Under any honest
enclosure an included zero forces $\ell-v\le0$; the displayed quadratic
control has perfect approximation but lacks theta arithmetic.

This exclusion leaves actual-height joint phase restrictions, correlations
among cores or with the anchor and correction, other groupings and
products, and contour or heat arguments open. Such an escape needs an
independently checked bound crossing the full paid threshold, with
coefficient variation and every cutoff included and an equally informed
direct comparator. No positive strip floor or Newman bound is acquired.

% 19 November 2026. Drafted for Edward Baker with OpenAI Codex assistance.
% E18 helper explicitly dispatched gpt-6.1-sol/ultra; backend unattested.
\section{An obstruction at actual dyadic phases}
\label{sec:actual-phase-recurrence}

The full-circle supremum of Section~\ref{sec:dyadic-group-obstruction}
can be replaced by the common dyadic phase at the actual height.
For the same independently relaxed core rotations, a paid slow-variation
argument still forces failure at some actual height in every length-19
interval within one natural chart. This is an existence statement
about the estimator, without a sign conclusion for the physical field.

All source quantities remain those of
\eqref{af:definitions}--\eqref{af:segment} and
\eqref{dg:groups}--\eqref{dg:matrix}. Thus $q=f-c$ retains the real-profile
value correction, $v$ retains both whole term-error sums and both
entire delta payments, and $D_\kappa$ remains a separate cost. Set
\begin{align}
 C_{\rm act}(x)&=\sum_rR_r(x,u(x)),\qquad u(x)=-\tau(x)\log2,\notag\\
 K_{\rm act}&=1-g-C_{\rm act},&
 \ell_{\rm act}&=\frac{K_{\rm act}}{|E|}-D_\kappa-|c|,\qquad
 m_{\rm act}=\ell_{\rm act}-v .
 \label{ap:allowance}
\end{align}
The radius still takes a supremum over each independent core rotation
$\eta_r$; only the common dyadic phase is now physical.
Reverse triangle in \eqref{dg:identity}, without taking a supremum
over $u$, gives $|EF|\ge K_{\rm act}$. Paying the complete errors proves
\begin{equation}
 \ell_{\rm act}\le|q|,\qquad
 m_E\le m_{\rm dyad}\le m_{\rm act}
       \le m_{\rm dir}=|f-c|-v.
 \label{ap:comparison}
\end{equation}
The divisor is the exact positive $|E|$, even for a negative numerator.
Direct evaluation receives the same coefficients, phases and payments.
Positive complete allowance would require exactly
$C_{\rm act}<1-g-|E|(D_\kappa+|c|+v)$.

\begin{theorem}[Scoped actual-phase recurrence]
\label{thm:actual-phase-recurrence}
On $\mathcal F_0$ of \eqref{af:segment}, every closed interval
$[a,a+19]$ contained in one fixed natural-cutoff chart contains an
actual height $x_*$ with
\begin{equation}
 C_{\rm act}(x_*)-(1-g(x_*))>\frac1{100},\qquad
 K_{\rm act}(x_*)<-\frac1{100},\qquad
 \ell_{\rm act}(x_*),\,m_{\rm act}(x_*)<-\frac1{800}.
 \label{ap:theorem}
\end{equation}
An excluded upper endpoint is interpreted in the old chart by its
continuous one-sided limit; the witness can be chosen strictly before
that endpoint. No failure at every height, at $x_0$ itself, or at a
seam is asserted.
\end{theorem}

\begin{proof}
Retain only the eight prime cores
$\mathcal P_*=\{7,11,13,17,19,23,29,31\}$, and write
$C_*(x,u)=\sum_{p\in\mathcal P_*}R_p(x,u)$.
The finite Fourier-mean argument of the preceding section and
\eqref{af:prime-mass} imply, at frozen height $a$,
\begin{equation}
 \langle C_*(a,\cdot)\rangle>
                  \frac{21}{20}(1-g(a)).
 \label{ap:mean}
\end{equation}
Indeed $\beta_p=b_p$, $a_p=p^{1/5}b_p$,
$h_p\ge(1-g)b_p$, and $\sigma<4/5$.

We first pay for coefficient variation. With $T=x/2$ and
$L_t=L_0+\alpha^2/40$, the explicit carrier derivatives are
\[
 \alpha'=-\frac1{2w^2}-\frac1{(w-1)^2}+\frac1{2w},\qquad
 \alpha''=\frac1{w^3}+\frac2{(w-1)^3}-\frac1{2w^2}.
\]
On $w=v+\ii T$, $2/5\le v\le3/5$ and $7000<T\le15000$,
$1000<|w|/(2\pi)<2500<e^8$ and $|\arg w|<\pi/2$ give
$|\alpha|<6$. Triangle then yields
\[
 |\alpha'|\le\frac1{2T}+\frac3{2T^2}
       <\frac{51}{100T},\qquad
 |\alpha''|<T^{-2},\qquad
 |L_t''|<\frac{51}{100T}+\frac7{20T^2}<\frac3{5T}.
\]
These use only $25/8<\pi<22/7$ and $e>8/3$; for example
$2500\cdot6561<8^8$.
Differentiating $s=s_0+\alpha(s_0)/20$ gives
$s'=-\ii/2-\ii\alpha'(s_0)/40$, hence
\begin{equation}
 |\sigma'|<\frac1{39x},\qquad
 |\tau'+1/2|<\frac1{39x},\qquad
 |g'|<\frac9{125x}<\frac1{10x}.
 \label{ap:velocities}
\end{equation}
For the last bound, conjugation of the fixed carrier logarithm gives
$\log g=\Re[L_t(2/5+\ii T)-L_t(3/5+\ii T)]$.
The horizontal integral of $L_t''$ has length $1/5$, and $dT/dx=1/2$,
so $|(\log g)'|<3/(25x)$; now use $g<3/5$.
The carrier logarithm keeps its fixed analytic branch; the cutoff
floor is not differentiated.

For these eight cores only $d=1,2$ occurs in the convolution.
The original indices at most48 are
\[
 7,14,28;\ 11,22,44;\ 13,26;\ 17,34;\
 19,38;\ 23,46;\ 29;\ 31.
\]
There are16, and each chart retains its $n\le N$ subset.
Each contributes $b_n$ at $m=n$ and $-b_2b_n$ at $m=2n$,
with the additional $n^{1/5}$ for $a_m$.
Using $\log n<4$, $\sigma>3/5$ and $n\ge7$ gives
$b_n n^{-\sigma}<2/5$,
$n^{1/5}b_n n^{-\sigma}<3/5$, and
$b_2 2^{-\sigma}<3/4$. The middle inequality follows from
$(5/3)^{10}<3^5<7^3$.
Consequently, summing over these eight cores only,
\begin{equation}
 B_{\rm abs}:=\sum_m|\beta_m|m^{-\sigma}<\frac{56}{5},
 \qquad
 A_{\rm abs}:=\sum_m|a_m|m^{-\sigma}<\frac{84}{5}.
 \label{ap:mass}
\end{equation}
These triangle bounds are used for variation, not for the physical
field. They include every edge through $m\le2N\le96$, and
$\log m<5$ since $96<(8/3)^5$.

Put $k=(1-g)/(1+g)$ and
$\mathsf Q_p=(1-g)p^{-\sigma}M_p$, so $R_p=\|\mathsf Q_p\|_2$.
Its $m=p2^j$ harmonic is
\[
 m^{-\sigma}
 \operatorname{diag}\bigl(k(\beta_m+ga_m),\beta_m-ga_m\bigr)
 \operatorname{Rot}(ju).
\]
Here $\operatorname{Rot}(v)$ is the real rotation matrix through $v$.
At fixed auxiliary $u$, its $g$-derivative has norm at most
$m^{-\sigma}(2|\beta_m|+|a_m|)$, because
$|k'|\le2$ and $|(kg)'|\le1$.
Its $\sigma$-derivative has norm at most
$\log m\,m^{-\sigma}(|\beta_m|+g|a_m|)$.
Since $\beta_m,a_m$ are constant on a fixed chart, summing these
derivative bounds with \eqref{ap:velocities}--\eqref{ap:mass} gives
\[
 \frac1x\left\{\frac5{39}\frac{532}{25}
                    +\frac1{10}\frac{196}{5}\right\}
 =\frac{6482}{975x}<\frac7x.
\]
Use $|\|A\|_2-\|B\|_2|\le\|A-B\|_2$ and integrate the smooth
matrix derivative, without assuming differentiability of its norm.
Uniformly in $u$ and $a\le x\le a+19$,
\begin{equation}
 |C_*(x,u)-C_*(a,u)|\le\frac{7(x-a)}a,\qquad
 |g(x)-g(a)|\le\frac{x-a}{10a}.
 \label{ap:variation}
\end{equation}
Both are strict for $x>a$ and equality zero at $x=a$.

Finally $\log2=2\int_0^{1/3}(1-r^2)^{-1}\,dr>2/3$ and
\eqref{ap:velocities} give
\[
 u(a+19)-u(a)>
 19\left(\frac12-\frac1{39\cdot14000}\right)\frac23
 =\frac{19}3-\frac{19}{819000}>\frac{44}7>2\pi.
\]
Since $u'>0$, there is a unique $b<a+19$ with
$u(b)=u(a)+2\pi$. The continuous periodic $C_*(a,\cdot)$
attains a maximum, which is at least its strictly bounded mean
\eqref{ap:mean}. Choose $x_*\in[a,b]$ attaining that phase.
With $d=x_*-a<19$, \eqref{ap:variation} proves
\[
 C_*(x_*,u(x_*))-(1-g(x_*))
 >\frac{1-g(a)}{20}-\frac{71d}{10a}
 >\frac1{50}-\frac{1349}{140000}
 =\frac{1451}{140000}>\frac1{100}.
\]
This also covers $d=0$. The full cost is at least $C_*$.
Thus $K_{\rm act}<-1/100$; the exact divisor $0<|E|<8$
gives $K_{\rm act}/|E|<-1/800$, after which the nonnegative
remaining payments prove \eqref{ap:theorem}.
\end{proof}

For the chart endpoints put $l_{33}=x_0$, $l_N=x_N$ for $N\ge34$,
$r_N=x_{N+1}$ for $N\le47$, and $r_{48}=30000$.
Every $r_N-l_N>38$: the first follows from
$(25/8)(4\cdot34^2-1/40)-14001>38$, interior lengths are
$4\pi(2N+1)>38$, and the last exceeds $7248/7$.
Thus the theorem covers the first and last length-19 interval of every
chart. Its witness lies strictly before an excluded old-chart upper
endpoint; the new chart owns the seam. The interval $[x_0,x_0+19]$
already prevents uniform positivity of this estimator on $\mathcal F_0$,
without determining its value at $x_0$. A fixed-reference transfer still
pays the signed jump and whole source payment of \eqref{uc:seams}.

The acquired phase-velocity and coefficient-variation information is
new relative to the auxiliary mean argument. Independent core rotations
remain relaxed, so no joint maximizing core phases are asserted to occur
physically. The exact direct margin's sign and measured evaluation cost
are open. Signed anchor/core aggregation, joint restrictions on the
$\eta_r$, other products, and joint correction or source payments remain
outside this exclusion.

The collision screen of Section~\ref{sec:arithmetic-obstruction}
applies separately to the complete literal source enclosure, the
finite-coefficient/phase-velocity estimates, and a strict complete
margin. A changed source can share coarse velocities or finite
coefficients while lacking the full theta enclosure; an included zero
always defeats the strict margin under an honest enclosure.
The exact quadratic control there has perfect approximation and lacks
theta arithmetic. Conditional and location qualifications of the other
controls remain unchanged. The rest of the two-dimensional finite
strip, effective tail onset and paid overlap, initial verification and
whole intermediate-time barrier remain unacquired. No positive strip
floor, physical zero, new coverage or Newman bound follows.

\begingroup
\raggedright
\footnotesize
\bibliographystyle{plainurl}
\bibliography{references}
\endgroup
\end{document}
